% =============================================================================
% ΚΕΦΑΛΑΙΟ 3: ΑΡΧΙΤΕΚΤΟΝΙΚΗ ΣΥΣΤΗΜΑΤΟΣ ΚΑΙ ΔΕΔΟΜΕΝΑ
% =============================================================================
% Αυτό το αρχείο περιέχει το πλήρες Κεφάλαιο 3 της διπλωματικής εργασίας
% "Μοντελοποίηση Αγορών Ηλεκτρικής Ενέργειας με Julia/JuMP"
% 
% Για ενσωμάτωση στο thesis.tex, χρησιμοποιήστε: % =============================================================================
% ΚΕΦΑΛΑΙΟ 3: ΑΡΧΙΤΕΚΤΟΝΙΚΗ ΣΥΣΤΗΜΑΤΟΣ ΚΑΙ ΔΕΔΟΜΕΝΑ
% =============================================================================
% Αυτό το αρχείο περιέχει το πλήρες Κεφάλαιο 3 της διπλωματικής εργασίας
% "Μοντελοποίηση Αγορών Ηλεκτρικής Ενέργειας με Julia/JuMP"
% 
% Για ενσωμάτωση στο thesis.tex, χρησιμοποιήστε: % =============================================================================
% ΚΕΦΑΛΑΙΟ 3: ΑΡΧΙΤΕΚΤΟΝΙΚΗ ΣΥΣΤΗΜΑΤΟΣ ΚΑΙ ΔΕΔΟΜΕΝΑ
% =============================================================================
% Αυτό το αρχείο περιέχει το πλήρες Κεφάλαιο 3 της διπλωματικής εργασίας
% "Μοντελοποίηση Αγορών Ηλεκτρικής Ενέργειας με Julia/JuMP"
% 
% Για ενσωμάτωση στο thesis.tex, χρησιμοποιήστε: % =============================================================================
% ΚΕΦΑΛΑΙΟ 3: ΑΡΧΙΤΕΚΤΟΝΙΚΗ ΣΥΣΤΗΜΑΤΟΣ ΚΑΙ ΔΕΔΟΜΕΝΑ
% =============================================================================
% Αυτό το αρχείο περιέχει το πλήρες Κεφάλαιο 3 της διπλωματικής εργασίας
% "Μοντελοποίηση Αγορών Ηλεκτρικής Ενέργειας με Julia/JuMP"
% 
% Για ενσωμάτωση στο thesis.tex, χρησιμοποιήστε: \input{chapter3_system_architecture}
% =============================================================================

\chapter{Αρχιτεκτονική Συστήματος και Δεδομένα}
\label{ch:system-architecture}

Στο παρόν κεφάλαιο περιγράφεται η αρχιτεκτονική του συστήματος που αναπτύχθηκε για την προσομοίωση της Προημερήσιας Αγοράς ηλεκτρικής ενέργειας. Η αρχιτεκτονική βασίζεται σε τρεις πυλώνες: την υποδομή δεδομένων που συλλέγει και αποθηκεύει δεδομένα από την πλατφόρμα διαφάνειας του ENTSO-E, τη βάση δεδομένων PostgreSQL που οργανώνει τα δεδομένα σε κατάλληλο σχήμα, και τον πυρήνα υπολογισμών σε Julia/JuMP που υλοποιεί τους αλγορίθμους βελτιστοποίησης.

Η γενική ροή δεδομένων του συστήματος απεικονίζεται στο Σχήμα~\ref{fig:system-architecture}. Τα δεδομένα συλλέγονται από την πλατφόρμα ENTSO-E μέσω ενός αυτοματοποιημένου pipeline σε Python, αποθηκεύονται στη βάση δεδομένων PostgreSQL, επεξεργάζονται από τους αλγορίθμους Julia/JuMP, και τα αποτελέσματα αποθηκεύονται πίσω στη βάση για ανάλυση και οπτικοποίηση.

% TODO: Εισαγωγή σχήματος αρχιτεκτονικής
\begin{figure}[htbp]
\centering
% \includegraphics[width=0.95\textwidth]{figures/system_architecture.pdf}
\fbox{\parbox{0.9\textwidth}{\centering\vspace{3cm}
\texttt{ENTSO-E Transparency Platform} \\
$\downarrow$ \textit{(Python pipeline)} \\
\texttt{PostgreSQL Database} \\
$\downarrow$ \textit{(Julia/JuMP)} \\
\texttt{Euphemia Engine} \\
$\downarrow$ \\
\texttt{simulations.energy\_prices}
\vspace{1cm}}}
\caption{Γενική αρχιτεκτονική του συστήματος προσομοίωσης αγοράς ηλεκτρικής ενέργειας.}
\label{fig:system-architecture}
\end{figure}


% =============================================================================
\section{Πηγές Δεδομένων}
\label{sec:data-sources}

Η κύρια πηγή δεδομένων για το σύστημα είναι η Πλατφόρμα Διαφάνειας του ENTSO-E (ENTSO-E Transparency Platform), η οποία παρέχει δημοσίως διαθέσιμα δεδομένα για τις ευρωπαϊκές αγορές ηλεκτρικής ενέργειας σύμφωνα με τον Κανονισμό (ΕΕ) 543/2013~\cite{entsoe-transparency, eu-regulation-543}.

\subsection{Πλατφόρμα Διαφάνειας ENTSO-E}
\label{subsec:entsoe-platform}

Το ENTSO-E (European Network of Transmission System Operators for Electricity) είναι ο οργανισμός που συντονίζει τους Διαχειριστές Συστημάτων Μεταφοράς (TSOs) της Ευρώπης. Η Πλατφόρμα Διαφάνειας που διατηρεί παρέχει πρόσβαση σε δεδομένα που καλύπτουν τη λειτουργία των ευρωπαϊκών αγορών ηλεκτρικής ενέργειας, συμπεριλαμβανομένων δεδομένων παραγωγής, κατανάλωσης, τιμών, διασυνοριακών ροών και διαθεσιμότητας μονάδων.

Η πλατφόρμα προσφέρει δύο κύριους τρόπους πρόσβασης. Ο πρώτος είναι το RESTful API, το οποίο επιτρέπει προγραμματιστική πρόσβαση στα δεδομένα μέσω HTTP αιτημάτων. Το API χρησιμοποιεί μορφή XML για τις απαντήσεις και απαιτεί κλειδί αυθεντικοποίησης (API key) που παρέχεται δωρεάν μετά από εγγραφή. Ο δεύτερος τρόπος είναι η SFTP πρόσβαση, η οποία παρέχει πρόσβαση σε προ-επεξεργασμένα αρχεία CSV ανά τύπο δεδομένων. Αυτή η μέθοδος είναι αποτελεσματικότερη για μαζική λήψη δεδομένων και χρησιμοποιείται στην παρούσα υλοποίηση.

\subsection{Τύποι Δεδομένων}
\label{subsec:data-types}

Για την προσομοίωση της Προημερήσιας Αγοράς, το σύστημα συλλέγει και επεξεργάζεται τους ακόλουθους τύπους δεδομένων:

\paragraph{Μονάδες Παραγωγής (Production and Generation Units)}
Ο πίνακας \texttt{production\_and\_generation\_units} περιέχει πληροφορίες για τις μονάδες παραγωγής ηλεκτρικής ενέργειας σε κάθε ζώνη προσφοράς. Τα δεδομένα περιλαμβάνουν τον κωδικό και το όνομα της μονάδας, τη ζώνη προσφοράς (bidding zone), τον τύπο καυσίμου (π.χ. φυσικό αέριο, λιγνίτης, υδροηλεκτρικά, αιολικά), την εγκατεστημένη ισχύ σε MW, και την κατάσταση λειτουργίας (commissioned/decommissioned). Αυτά τα δεδομένα χρησιμοποιούνται για τη δημιουργία του συνόλου των διαθέσιμων μονάδων στο πρόβλημα Unit Commitment.

\paragraph{Πραγματικό Φορτίο (Actual Total Load)}
Ο πίνακας \texttt{actual\_total\_load} περιέχει ιστορικά δεδομένα κατανάλωσης ηλεκτρικής ενέργειας ανά ζώνη προσφοράς με ωριαία ή τεταρτοωριαία ανάλυση. Η ζήτηση αποτελεί βασική παράμετρο εισόδου τόσο στο πρόβλημα Unit Commitment όσο και στην εκκαθάριση της αγοράς.

\paragraph{Προβλέψεις Παραγωγής ΑΠΕ (Generation Forecasts for Wind and Solar)}
Ο πίνακας \texttt{generation\_forecasts\_for\_wind\_and\_solar} περιέχει προβλέψεις παραγωγής από αιολικά και φωτοβολταϊκά συστήματα. Δεδομένου ότι η παραγωγή από ΑΠΕ θεωρείται μη κατανεμόμενη (non-dispatchable), οι προβλέψεις αφαιρούνται από τη συνολική ζήτηση για τον υπολογισμό της καθαρής ζήτησης (net demand) που πρέπει να καλυφθεί από τις θερμικές μονάδες.

\paragraph{Διαθέσιμη Ικανότητα Μεταφοράς (Offered Transfer Capacities)}
Ο πίνακας \texttt{offered\_transfer\_capacities\_implicit} περιέχει δεδομένα για τη διαθέσιμη ικανότητα μεταφοράς (ATC) μεταξύ γειτονικών ζωνών προσφοράς. Τα δεδομένα περιλαμβάνουν τη ζώνη προέλευσης και προορισμού, τη χρονική περίοδο, και τη διαθέσιμη ικανότητα σε MW. Αυτά τα δεδομένα είναι απαραίτητα για την πολυζωνική εκκαθάριση της αγοράς με περιορισμούς δικτύου.

\paragraph{Μη Διαθεσιμότητα Μονάδων (Unavailability of Production Units)}
Ο πίνακας \texttt{unavailability\_of\_production\_and\_generation\_units} περιέχει πληροφορίες για προγραμματισμένες ή απρόβλεπτες διακοπές λειτουργίας μονάδων παραγωγής. Κάθε εγγραφή περιλαμβάνει τον κωδικό της μονάδας, την περίοδο μη διαθεσιμότητας, τον τύπο (προγραμματισμένη ή έκτακτη), την κατάσταση (ενεργή, ακυρωμένη), και την εναπομένουσα διαθέσιμη ισχύ. Αυτά τα δεδομένα χρησιμοποιούνται για τον αποκλεισμό ή τη μείωση της διαθέσιμης ισχύος μονάδων κατά την επίλυση του Unit Commitment.

\paragraph{Πραγματική Παραγωγή ανά Μονάδα (Actual Generation Output)}
Ο πίνακας \texttt{actual\_generation\_output\_per\_generation\_unit} περιέχει ιστορικά δεδομένα πραγματικής παραγωγής ανά μονάδα, με χρονική ανάλυση 15 λεπτών ή 1 ώρας. Κάθε εγγραφή περιλαμβάνει τον κωδικό της μονάδας, τη χρονοσφραγίδα (UTC), τον κωδικό ανάλυσης (\texttt{PT15M}, \texttt{PT60M}) και τη μετρούμενη παραγωγή σε MW. Τα δεδομένα αυτά αξιοποιούνται για τη στατιστική εξαγωγή τεχνικών παραμέτρων (ρυθμοί μεταβολής, ελάχιστη ισχύς, χρόνοι λειτουργίας/αδράνειας) μέσω ανάλυσης ιστορικού παραθύρου 3 μηνών.

\paragraph{Τιμές Ενέργειας (Energy Prices)}
Ο πίνακας \texttt{energy\_prices} περιέχει τις πραγματικές τιμές εκκαθάρισης της Προημερήσιας Αγοράς ανά ζώνη προσφοράς. Αυτά τα δεδομένα χρησιμοποιούνται για την επικύρωση (validation) των αποτελεσμάτων του συστήματος προσομοίωσης.

\subsection{Διαθεσιμότητα και Ποιότητα Δεδομένων}
\label{subsec:data-quality}

Η ποιότητα και η διαθεσιμότητα των δεδομένων του ENTSO-E ποικίλλει ανάλογα με τη ζώνη προσφοράς και τον τύπο δεδομένων. Ορισμένες ζώνες (π.χ. Γερμανία, Γαλλία, Σκανδιναβικές χώρες) διαθέτουν πλήρη και υψηλής ποιότητας δεδομένα, ενώ άλλες μπορεί να έχουν ελλείψεις ή ασυνέπειες. Για την Ελλάδα (ζώνη GR), τα δεδομένα είναι γενικά διαθέσιμα από το 2015 και μετά με καλή ποιότητα.

Οι κύριες προκλήσεις που αντιμετωπίζονται στην επεξεργασία των δεδομένων περιλαμβάνουν διαφορετικές χρονικές αναλύσεις (ωριαία, ημίωρη, τεταρτοωριαία) που απαιτούν ομογενοποίηση, ελλιπείς εγγραφές που απαιτούν συμπλήρωση ή αποκλεισμό, διαφορετικές ζώνες ώρας και μετάβαση θερινής/χειμερινής ώρας, και αλλαγές στη δομή των δεδομένων κατά τη διάρκεια του χρόνου.


% =============================================================================
\section{Αγωγός Δεδομένων (Data Pipeline)}
\label{sec:data-pipeline}

Για τη συστηματική συλλογή και ενημέρωση των δεδομένων από την πλατφόρμα ENTSO-E, αναπτύχθηκε ένας αυτοματοποιημένος αγωγός δεδομένων (data pipeline) σε Python. Ο αγωγός υλοποιεί ένα σύστημα σταδιακών ενημερώσεων (incremental updates) που επιτρέπει την αποτελεσματική ενημέρωση της βάσης δεδομένων χωρίς την ανάγκη πλήρους επαναφόρτωσης~\cite{kleppmann2017}.

\subsection{Αρχιτεκτονική Pipeline}
\label{subsec:pipeline-architecture}

Ο αγωγός δεδομένων αποτελείται από τα ακόλουθα βασικά στοιχεία:

\paragraph{Κεντρική Βιβλιοθήκη (helper\_functions.py)}
Περιέχει τις βασικές λειτουργίες που χρησιμοποιούνται από όλα τα scripts ενημέρωσης, συμπεριλαμβανομένων συναρτήσεων αυθεντικοποίησης με την πλατφόρμα ENTSO-E, ανακάλυψης και φιλτραρίσματος αρχείων, λήψης και επεξεργασίας CSV, μαζικής φόρτωσης στη βάση δεδομένων, και διαχείρισης αρχείων καταγραφής.

\paragraph{Scripts Ενημέρωσης (update\_*.py)}
Για κάθε τύπο δεδομένων υπάρχει ένα ξεχωριστό script ενημέρωσης που καθορίζει τις παραμέτρους ειδικά για αυτόν τον τύπο (όνομα πίνακα, φάκελος ENTSO-E, διαχωριστικό αρχείου, κωδικοποίηση κ.λπ.). Αυτή η αρχιτεκτονική επιτρέπει την εύκολη προσθήκη νέων τύπων δεδομένων και την ανεξάρτητη εκτέλεση κάθε ενημέρωσης.

\paragraph{Εργαλείο Δημιουργίας Scripts (generate\_incremental\_scripts.py)}
Αυτοματοποιεί τη δημιουργία των scripts ενημέρωσης με βάση πρότυπα, διασφαλίζοντας συνέπεια στη δομή και τη λειτουργικότητα.

\subsection{Σύστημα Σταδιακών Ενημερώσεων}
\label{subsec:incremental-updates}

Η κεντρική καινοτομία του αγωγού δεδομένων είναι το σύστημα σταδιακών ενημερώσεων που αποφεύγει την επαναφόρτωση ολόκληρου του συνόλου δεδομένων σε κάθε εκτέλεση. Η ροή εργασιών περιγράφεται ως εξής:

Αρχικά, το σύστημα εκτελεί αυθεντικοποίηση με την πλατφόρμα ENTSO-E χρησιμοποιώντας διαπιστευτήρια που αποθηκεύονται σε μεταβλητές περιβάλλοντος. Στη συνέχεια, συνδέεται μέσω SFTP και ανακαλύπτει τα διαθέσιμα αρχεία στον αντίστοιχο φάκελο, καταγράφοντας το όνομα και την ημερομηνία τελευταίας τροποποίησης κάθε αρχείου. Το σύστημα συγκρίνει τη λίστα με τον πίνακα \texttt{file\_update\_log} της βάσης δεδομένων και εντοπίζει τα νέα ή τροποποιημένα αρχεία. Τα αρχεία επεξεργάζονται σε παρτίδες (batches) καθορισμένου μεγέθους (προεπιλογή: 50 αρχεία) για αποτελεσματική χρήση μνήμης. Κάθε αρχείο CSV αναλύεται, καθαρίζεται (data cleaning) και μετασχηματίζεται στη μορφή του πίνακα προορισμού. Τέλος, τα δεδομένα φορτώνονται στη βάση χρησιμοποιώντας τη μέθοδο PostgreSQL COPY για μέγιστη απόδοση, και ο πίνακας \texttt{file\_update\_log} ενημερώνεται με τα στοιχεία των επεξεργασμένων αρχείων.

\subsection{Βελτιστοποίηση Απόδοσης}
\label{subsec:pipeline-performance}

Για τη βελτιστοποίηση της απόδοσης του αγωγού δεδομένων, υλοποιήθηκαν διάφορες τεχνικές. Η χρήση της μεθόδου PostgreSQL COPY αντί για πολλαπλές εντολές INSERT προσφέρει σημαντικά υψηλότερη απόδοση για μαζική φόρτωση δεδομένων. Το ρυθμιζόμενο μέγεθος παρτίδας (batch size) επιτρέπει την εξισορρόπηση μεταξύ χρήσης μνήμης και απόδοσης. Η επεξεργασία σε παρτίδες αντί της φόρτωσης ολόκληρου του αρχείου στη μνήμη επιτρέπει τη διαχείριση πολύ μεγάλων συνόλων δεδομένων. Ο μηχανισμός παρακολούθησης αρχείων μέσω του \texttt{file\_update\_log} αποφεύγει την περιττή επαναεπεξεργασία αρχείων.

\subsection{Χρονοπρογραμματισμός}
\label{subsec:scheduling}

Η αυτοματοποιημένη εκτέλεση του αγωγού δεδομένων επιτυγχάνεται μέσω GitHub Actions, με τις εργασίες να εκτελούνται σε καθημερινή βάση. Η ροή εργασιών περιλαμβάνει την εκτέλεση του \texttt{generate\_incremental\_scripts.py} για τη δημιουργία των τελευταίων scripts, ακολουθούμενη από την εκτέλεση κάθε script ενημέρωσης για τους διάφορους τύπους δεδομένων.


% =============================================================================
\section{Σχήμα Βάσης Δεδομένων}
\label{sec:database-schema}

Η βάση δεδομένων PostgreSQL οργανώνεται σε δύο κύρια σχήματα (schemas): το σχήμα \texttt{entsoe} που περιέχει τα ακατέργαστα δεδομένα από την πλατφόρμα ENTSO-E, και το σχήμα \texttt{simulations} που περιέχει τα αποτελέσματα των προσομοιώσεων.

\subsection{Σχήμα ENTSO-E}
\label{subsec:entsoe-schema}

Το σχήμα \texttt{entsoe} περιέχει τους πίνακες που αποθηκεύουν τα δεδομένα από την πλατφόρμα διαφάνειας. Οι κύριοι πίνακες παρουσιάζονται στον Πίνακα~\ref{tab:entsoe-tables}.

\begin{table}[htbp]
\centering
\caption{Κύριοι πίνακες του σχήματος \texttt{entsoe}}
\label{tab:entsoe-tables}
\begin{tabular}{@{}p{5.5cm}p{8cm}@{}}
\toprule
\textbf{Πίνακας} & \textbf{Περιγραφή} \\
\midrule
\texttt{production\_and\_generation\_units} & Στοιχεία μονάδων παραγωγής \\
\texttt{actual\_total\_load} & Ιστορικά δεδομένα κατανάλωσης \\
\texttt{actual\_generation\_output\_per\_generation\_unit} & Ιστορική παραγωγή ανά μονάδα \\
\texttt{generation\_forecasts\_for\_wind\_and\_solar} & Προβλέψεις παραγωγής ΑΠΕ \\
\texttt{offered\_transfer\_capacities\_implicit} & Διαθέσιμη ικανότητα μεταφοράς \\
\texttt{unavailability\_of\_production\_and\_generation\_units} & Μη διαθεσιμότητα μονάδων \\
\texttt{energy\_prices} & Τιμές εκκαθάρισης αγοράς \\
\texttt{physical\_flows} & Φυσικές ροές διασυνδέσεων \\
\texttt{file\_update\_log} & Καταγραφή επεξεργασμένων αρχείων \\
\bottomrule
\end{tabular}
\end{table}

\paragraph{Πίνακας production\_and\_generation\_units}
Ο πίνακας αυτός περιέχει τα στοιχεία των μονάδων παραγωγής με τα ακόλουθα βασικά πεδία: \texttt{generation\_unit\_code} (μοναδικός κωδικός μονάδας), \texttt{generation\_unit\_name} (όνομα μονάδας), \texttt{area\_map\_code} (κωδικός ζώνης προσφοράς, π.χ. ``GR''), \texttt{production\_unit\_type} (τύπος καυσίμου), \texttt{generation\_unit\_installed\_capacity\_mw} (εγκατεστημένη ισχύς), \texttt{production\_unit\_status} (κατάσταση λειτουργίας), \texttt{valid\_from} και \texttt{valid\_to} (περίοδος ισχύος των στοιχείων).

Η ερώτηση για την ανάκτηση διαθέσιμων μονάδων για μια συγκεκριμένη ζώνη και ημερομηνία φιλτράρει τις μονάδες με βάση την κατάσταση \texttt{COMMISSIONED}, τον τύπο περιοχής \texttt{BZN} ή \texttt{BZN/CTA}, και την περίοδο ισχύος.

\paragraph{Πίνακας actual\_total\_load}
Ο πίνακας αυτός περιέχει χρονοσειρές κατανάλωσης με τα πεδία: \texttt{datetime\_utc} (χρονοσφραγίδα σε UTC), \texttt{area\_map\_code} (ζώνη προσφοράς), \texttt{total\_load\_mw} (κατανάλωση σε MW), και \texttt{resolution} (χρονική ανάλυση).

\paragraph{Πίνακας actual\_generation\_output\_per\_generation\_unit}
Ο πίνακας αυτός περιέχει ιστορικά δεδομένα πραγματικής παραγωγής ανά μονάδα με τα πεδία: \texttt{generation\_unit\_code} (κωδικός μονάδας), \texttt{datetime\_utc} (χρονοσφραγίδα σε UTC), \texttt{actual\_generation\_output\_mw} (μετρούμενη παραγωγή σε MW), και \texttt{resolution} (χρονική ανάλυση, συνήθως 15 λεπτά ή 1 ώρα). Τα δεδομένα αυτά χρησιμοποιούνται για την εξαγωγή τεχνικών παραμέτρων (ρυθμοί μεταβολής, ελάχιστη ισχύς, χρόνοι λειτουργίας/αδράνειας) μέσω στατιστικής ανάλυσης, όπως περιγράφεται στην Ενότητα~\ref{subsec:parameter-inference}.

\paragraph{Πίνακας unavailability\_of\_production\_and\_generation\_units}
Ο πίνακας αυτός περιέχει πληροφορίες διακοπών με τα πεδία: \texttt{asset\_code} (αντιστοιχεί στο \texttt{generation\_unit\_code}), \texttt{start\_outage\_utc} και \texttt{end\_outage\_utc} (περίοδος διακοπής), \texttt{type} (``Planned'' ή ``Forced''), \texttt{status} (``Active'', ``Cancelled'', ``Withdrawn''), και \texttt{available\_capacity\_mw} (εναπομένουσα διαθέσιμη ισχύς).

\subsection{Σχήμα Προσομοιώσεων}
\label{subsec:simulations-schema}

Το σχήμα \texttt{simulations} περιέχει τους πίνακες για την αποθήκευση των αποτελεσμάτων των προσομοιώσεων. Οι κύριοι πίνακες παρουσιάζονται στον Πίνακα~\ref{tab:simulations-tables}.

\begin{table}[htbp]
\centering
\caption{Κύριοι πίνακες του σχήματος \texttt{simulations}}
\label{tab:simulations-tables}
\begin{tabular}{@{}p{5cm}p{8cm}@{}}
\toprule
\textbf{Πίνακας} & \textbf{Περιγραφή} \\
\midrule
\texttt{energy\_prices} & Υπολογισμένες τιμές εκκαθάρισης \\
\texttt{optimization\_runs} & Μεταδεδομένα εκτελέσεων βελτιστοποίησης \\
\texttt{transmission\_flows} & Διασυνοριακές ροές ισχύος \\
\texttt{generator\_inferred\_parameters} & Εξαχθείσες τεχνικές παράμετροι μονάδων \\
\texttt{uc\_results} & Αποτελέσματα Unit Commitment (σύνοψη) \\
\texttt{uc\_generation} & Παραγωγή/δέσμευση ανά μονάδα και περίοδο \\
\texttt{uc\_net\_demand} & Καθαρή ζήτηση ανά περίοδο \\
\bottomrule
\end{tabular}
\end{table}

\paragraph{Πίνακας energy\_prices}
Ο πίνακας αυτός αποθηκεύει τις υπολογισμένες τιμές εκκαθάρισης με τα πεδία: \texttt{bidding\_zone} (ζώνη προσφοράς), \texttt{date} (ημερομηνία), \texttt{hour} (ώρα 1-24), \texttt{price} (τιμή σε €/MWh), \texttt{created\_at} (χρονοσφραγίδα δημιουργίας), και \texttt{run\_id} (αναφορά στην εκτέλεση βελτιστοποίησης).

\paragraph{Πίνακας transmission\_flows}
Ο πίνακας αυτός αποθηκεύει τις υπολογισμένες διασυνοριακές ροές με τα πεδία: \texttt{source\_zone} και \texttt{target\_zone} (ζώνες προέλευσης και προορισμού), \texttt{date} (ημερομηνία), \texttt{hour} (ώρα), \texttt{flow\_mw} (ροή σε MW), και \texttt{run\_id}.

\paragraph{Πίνακας generator\_inferred\_parameters}
Ο πίνακας αυτός λειτουργεί ως κρυφή μνήμη (cache) για τις τεχνικές παραμέτρους μονάδων παραγωγής που εξάγονται από ιστορικά δεδομένα. Τα βασικά πεδία περιλαμβάνουν: \texttt{generator\_code} (κωδικός μονάδας), \texttt{ramp\_up} και \texttt{ramp\_down} (ρυθμοί μεταβολής ως κλάσμα της $P^{\max}$ ανά ώρα), \texttt{p\_min} (ελάχιστη ισχύς ως κλάσμα της $P^{\max}$), \texttt{min\_uptime} και \texttt{min\_downtime} (ελάχιστοι χρόνοι λειτουργίας/αδράνειας σε ώρες), και \texttt{inference\_date} (ημερομηνία εξαγωγής για έλεγχο λήξης). Η χρήση του πίνακα αυτού μειώνει τον χρόνο ανάκτησης παραμέτρων από περίπου 17 λεπτά σε 2 δευτερόλεπτα.

\paragraph{Πίνακες uc\_results, uc\_generation, uc\_net\_demand}
Οι τρεις αυτοί πίνακες αποτελούν τον μηχανισμό κρυφής μνήμης (caching) για τα αποτελέσματα του επιλυτή Unit Commitment. Ο πίνακας \texttt{uc\_results} αποθηκεύει τη σύνοψη κάθε εκτέλεσης με κλειδί \texttt{(bidding\_zone, market\_date, code\_version)} και περιλαμβάνει τα πεδία: \texttt{status} (κατάσταση τερματισμού), \texttt{solver} (επιλυτής), \texttt{resolution\_minutes} (χρονική ανάλυση), \texttt{total\_cost}, \texttt{production\_cost}, \texttt{startup\_cost}, \texttt{noload\_cost} (συνιστώσες κόστους), και αριθμούς εκκινήσεων (\texttt{hot/warm/cold\_startups}). Ο πίνακας \texttt{uc\_generation} αποθηκεύει τα αναλυτικά αποτελέσματα παραγωγής, δέσμευσης και εκκίνησης ανά μονάδα και περίοδο, ενώ ο πίνακας \texttt{uc\_net\_demand} αποθηκεύει την καθαρή ζήτηση και την παραγωγή ΑΠΕ ανά περίοδο. Η αποθήκευση ανέρχεται σε περίπου 195 KB ανά ζώνη/ημέρα.


% =============================================================================
\section{Οικοσύστημα Julia/JuMP}
\label{sec:julia-jump}

Για την υλοποίηση των αλγορίθμων βελτιστοποίησης επιλέχθηκε η γλώσσα προγραμματισμού Julia σε συνδυασμό με το πλαίσιο μοντελοποίησης JuMP~\cite{bezanson2017, jump2017}. Αυτή η επιλογή βασίζεται σε τεχνικά πλεονεκτήματα που καθιστούν το συγκεκριμένο οικοσύστημα ιδανικό για επιστημονικούς υπολογισμούς και προβλήματα βελτιστοποίησης.

\subsection{Γλώσσα Προγραμματισμού Julia}
\label{subsec:julia-language}

Η Julia είναι μια σύγχρονη γλώσσα προγραμματισμού υψηλού επιπέδου σχεδιασμένη για επιστημονικούς υπολογισμούς. Τα κύρια χαρακτηριστικά της που την καθιστούν κατάλληλη για την παρούσα εφαρμογή είναι τα ακόλουθα.

Η Julia επιτυγχάνει υψηλή απόδοση μέσω της μεταγλώττισης just-in-time (JIT) με χρήση του LLVM compiler framework. Ο κώδικας Julia μεταγλωττίζεται σε εγγενή κώδικα μηχανής, επιτυγχάνοντας ταχύτητες συγκρίσιμες με C και Fortran, χωρίς την ανάγκη για χειροκίνητη βελτιστοποίηση ή κλήση εξωτερικών βιβλιοθηκών.

Παρά την υψηλή απόδοση, η σύνταξη της Julia είναι εκφραστική και φιλική στον χρήστη, παρόμοια με Python ή MATLAB. Αυτό επιτρέπει την ταχεία ανάπτυξη πρωτοτύπων και την εύκολη ανάγνωση του κώδικα.

Η Julia διαθέτει ένα εξελιγμένο σύστημα τύπων με πολλαπλή αποστολή (multiple dispatch) που επιτρέπει τη δημιουργία γενικευμένου κώδικα που είναι ταυτόχρονα αποτελεσματικός και επαναχρησιμοποιήσιμος.

Η Julia παρέχει ενσωματωμένη υποστήριξη για παραλληλισμό και κατανεμημένους υπολογισμούς, διευκολύνοντας την κλιμάκωση σε πολλαπλούς πυρήνες ή μηχανές.

\subsection{Πλαίσιο Μοντελοποίησης JuMP}
\label{subsec:jump-framework}

Το JuMP (Julia for Mathematical Programming) είναι ένα domain-specific language (DSL) ενσωματωμένο στη Julia για τη διατύπωση και επίλυση προβλημάτων μαθηματικού προγραμματισμού. Αποτελεί ένα από τα πιο δημοφιλή και ώριμα πλαίσια μοντελοποίησης βελτιστοποίησης στον ακαδημαϊκό και βιομηχανικό χώρο~\cite{jump2017}.

Τα πλεονεκτήματα του JuMP περιλαμβάνουν μια αλγεβρική μοντελοποίηση που επιτρέπει τη διατύπωση προβλημάτων βελτιστοποίησης με τρόπο που μοιάζει με τη μαθηματική τους διατύπωση. Η σύνταξη για μεταβλητές, περιορισμούς και αντικειμενικές συναρτήσεις είναι διαισθητική και εκφραστική. Το JuMP υποστηρίζει πολλαπλούς επιλυτές (solvers) μέσω μιας ενιαίας διεπαφής, επιτρέποντας την εύκολη εναλλαγή μεταξύ επιλυτών χωρίς αλλαγές στον κώδικα μοντελοποίησης. Η απόδοση του JuMP είναι συγκρίσιμη ή καλύτερη από εμπορικά εργαλεία μοντελοποίησης όπως το AMPL και το GAMS. Τέλος, το JuMP είναι εντελώς ελεύθερο και ανοιχτού κώδικα, διατίθεται υπό την άδεια MIT.

Ένα τυπικό παράδειγμα χρήσης του JuMP για τη δημιουργία ενός μοντέλου βελτιστοποίησης φαίνεται στο ακόλουθο απόσπασμα κώδικα:

\begin{samepage}
\begin{minted}[fontsize=\small,linenos,frame=lines]{julia}
using JuMP, HiGHS

# Δημιουργία μοντέλου με τον επιλυτή HiGHS
model = Model(HiGHS.Optimizer)

# Ορισμός μεταβλητών
@variable(model, x >= 0)
@variable(model, y >= 0)

# Ορισμός περιορισμών
@constraint(model, x + 2y <= 10)
@constraint(model, 3x + y <= 15)

# Ορισμός αντικειμενικής συνάρτησης
@objective(model, Max, 2x + 3y)

# Επίλυση
optimize!(model)
\end{minted}
\end{samepage}

\subsection{Επιλυτής HiGHS}
\label{subsec:highs-solver}

Ως κύριος επιλυτής για τα προβλήματα βελτιστοποίησης χρησιμοποιείται ο HiGHS, ένας σύγχρονος επιλυτής ανοιχτού κώδικα για προβλήματα γραμμικού προγραμματισμού (LP), μικτού ακέραιου προγραμματισμού (MIP) και τετραγωνικού προγραμματισμού (QP)~\cite{highs2024}.

Ο HiGHS αναπτύσσεται στο Πανεπιστήμιο του Εδιμβούργου και αποτελεί έναν από τους ταχύτερους ανοιχτού κώδικα επιλυτές. Περιλαμβάνει υλοποιήσεις του αλγορίθμου dual revised simplex με εκτεταμένη παραλληλοποίηση, μεθόδων εσωτερικού σημείου (interior point methods), και αλγορίθμων branch-and-bound για MIP με cutting planes και heuristics.

Η διεπαφή με τη Julia παρέχεται μέσω του πακέτου HiGHS.jl, το οποίο ενσωματώνεται απρόσκοπτα με το JuMP. Η επιλογή του HiGHS ως προεπιλεγμένου επιλυτή βασίζεται στην άδεια ανοιχτού κώδικα (MIT) που επιτρέπει ελεύθερη χρήση, στην υψηλή απόδοση συγκρίσιμη με εμπορικούς επιλυτές, στην ενεργή ανάπτυξη και υποστήριξη, και στη συμβατότητα με όλα τα απαιτούμενα είδη προβλημάτων (LP, MILP).

Το σύστημα υποστηρίζει επίσης εμπορικούς επιλυτές όπως ο Gurobi και ο CPLEX για περιπτώσεις που απαιτείται ακόμη υψηλότερη απόδοση, με την εναλλαγή να γίνεται μέσω μιας απλής παραμέτρου.

\subsection{Βασικά Πακέτα Julia}
\label{subsec:julia-packages}

Εκτός από το JuMP και τον HiGHS, το σύστημα χρησιμοποιεί μια σειρά από πακέτα Julia για διάφορες λειτουργίες. Το πακέτο DataFrames.jl χρησιμοποιείται για τη διαχείριση και επεξεργασία δεδομένων σε μορφή πινάκων, με λειτουργικότητα παρόμοια με τα pandas της Python. Το πακέτο LibPQ.jl παρέχει σύνδεση με τη βάση δεδομένων PostgreSQL μέσω της βιβλιοθήκης libpq, υποστηρίζοντας παραμετροποιημένα ερωτήματα και διαχείριση συνδέσεων. Το πακέτο CSV.jl χρησιμοποιείται για την ανάγνωση και εγγραφή αρχείων CSV με υψηλή απόδοση. Το πακέτο Dates.jl, που αποτελεί μέρος της standard library, χρησιμοποιείται για τη διαχείριση ημερομηνιών και χρόνων. Το πακέτο Statistics.jl, επίσης μέρος της standard library, χρησιμοποιείται για στατιστικούς υπολογισμούς---κυρίως τη συνάρτηση \texttt{quantile} για τον υπολογισμό εκατοστημορίων κατά την εξαγωγή τεχνικών παραμέτρων μονάδων από ιστορικά δεδομένα. Τέλος, το πακέτο DotEnv.jl χρησιμοποιείται για τη φόρτωση μεταβλητών περιβάλλοντος από αρχεία \texttt{.env} για τη διαχείριση διαπιστευτηρίων.


% =============================================================================
\section{Αρχιτεκτονική Εφαρμογής}
\label{sec:application-architecture}

Η εφαρμογή Julia οργανώνεται ως ένα ενιαίο module ονόματι \texttt{Euphemia} που περιέχει όλη τη λειτουργικότητα του συστήματος. Η δομή του module ακολουθεί τις βέλτιστες πρακτικές της Julia για οργάνωση κώδικα και επαναχρησιμοποίηση~\cite{julia-docs}.

\subsection{Δομή Module}
\label{subsec:module-structure}

Το module \texttt{Euphemia} αποτελείται από τα ακόλουθα υπο-modules και αρχεία, όπως φαίνεται στο Σχήμα~\ref{fig:module-structure}:

\begin{figure}[htbp]
\centering
\begin{verbatim}
Euphemia.jl (Main Module)
├── dbutils.jl                 # Database utilities
├── Generators.jl              # Generator data model
├── Loads.jl                   # Load data model
├── Renewables.jl              # RES forecasts
├── FuelTypeParameters.jl      # Fuel-specific parameters
├── MarketOrders.jl            # Order types
├── UnitCommitment.jl          # UC solver
├── BiddingStrategy.jl         # UC → Market orders
├── Network.jl                 # ATC constraints
├── MPCC.jl                    # Market clearing
└── AlternativeOrderBook.jl    # Fast order generation
\end{verbatim}
\caption{Δομή του module \texttt{Euphemia}.}
\label{fig:module-structure}
\end{figure}

\paragraph{Κεντρικό Module (Euphemia.jl)}
Αποτελεί το σημείο εισόδου της εφαρμογής. Διαχειρίζεται τις εξαρτήσεις, συντονίζει τα υπο-modules, και εξάγει τη δημόσια διεπαφή (public API). Περιλαμβάνει επίσης τη λογική επιλογής επιλυτή και την κεντρική ροή εργασιών.

\paragraph{Βοηθητικά Βάσης Δεδομένων (dbutils.jl)}
Περιέχει συναρτήσεις για τη σύνδεση με τη βάση δεδομένων PostgreSQL και την εκτέλεση ερωτημάτων. Υλοποιεί το pattern \texttt{withdb} για ασφαλή διαχείριση συνδέσεων με αυτόματο κλείσιμο.

\paragraph{Μοντέλα Δεδομένων (Generators.jl, Loads.jl, Renewables.jl)}
Ορίζουν τις δομές δεδομένων για τις οντότητες του συστήματος και παρέχουν συναρτήσεις για την ανάκτησή τους από τη βάση δεδομένων. Η δομή \texttt{Generator} περιλαμβάνει πεδία για τον κωδικό, το όνομα, τη ζώνη προσφοράς, τον τύπο καυσίμου, τα όρια ισχύος, το οριακό κόστος, και προαιρετικούς τεχνικούς περιορισμούς (ρυθμοί μεταβολής, ελάχιστοι χρόνοι λειτουργίας/αδράνειας) που εξάγονται αυτόματα από ιστορικά δεδομένα. Η δομή \texttt{InitialConditions} αναπαριστά την αρχική κατάσταση κάθε μονάδας στη χρονική στιγμή $t=0$ (δέσμευση, ισχύς εξόδου, θερμοκρασιακή κατάσταση) και εξάγεται από ιστορικά δεδομένα 72 ωρών μέσω μαζικού ερωτήματος. Το αρχείο \texttt{Generators.jl} περιλαμβάνει τις συναρτήσεις εξαγωγής παραμέτρων (\texttt{infer\_ramp\_rates}, \texttt{infer\_p\_min}, \texttt{infer\_uptime\_downtime}), τη συνάρτηση \texttt{get\_initial\_conditions} για αυτόματο προσδιορισμό αρχικών συνθηκών, και τον μηχανισμό αποθήκευσης σε κρυφή μνήμη βάσης δεδομένων.

\paragraph{Τύποι Εντολών (MarketOrders.jl)}
Ορίζει τις δομές δεδομένων για τους διάφορους τύπους εντολών αγοράς, συμπεριλαμβανομένων των \texttt{SimpleOrder}, \texttt{BlockOrder}, \texttt{LinkedBlockOrder} κ.λπ.

\paragraph{Unit Commitment (UnitCommitment.jl)}
Υλοποιεί τον επιλυτή για το πρόβλημα Unit Commitment χρησιμοποιώντας JuMP/HiGHS. Περιλαμβάνει τη διατύπωση του μοντέλου με τριμερή αντικειμενική συνάρτηση (κόστος παραγωγής, κόστος εκκίνησης ανά θερμοκρασιακή κατάσταση, κόστος αδράνειας), την ενσωμάτωση αρχικών συνθηκών, ρυθμιζόμενες παραμέτρους επιλυτή (MIP gap, χρονικό όριο μέσω MOI), καθώς και μηχανισμό αποθήκευσης αποτελεσμάτων σε κρυφή μνήμη βάσης δεδομένων.

\paragraph{Στρατηγική Προσφορών (BiddingStrategy.jl)}
Υλοποιεί τη μετατροπή των αποτελεσμάτων του Unit Commitment σε εντολές αγοράς σύμφωνα με τις τρεις στρατηγικές που περιγράφηκαν στο Κεφάλαιο~\ref{sec:bidding-strategies}.

\paragraph{Δίκτυο (Network.jl)}
Διαχειρίζεται τους περιορισμούς δικτύου και τη διαθέσιμη ικανότητα μεταφοράς (ATC) μεταξύ ζωνών προσφοράς. Παρέχει λειτουργίες για την ανάκτηση δεδομένων ATC από το ENTSO-E και την προσθήκη περιορισμών στο μοντέλο βελτιστοποίησης.

\paragraph{Εκκαθάριση Αγοράς (MPCC.jl)}
Υλοποιεί τον αλγόριθμο εκκαθάρισης της αγοράς με τη μέθοδο MPCC. Περιλαμβάνει τη δομή \texttt{MPCCOrderBook} για τη διαχείριση των εντολών και τη συνάρτηση \texttt{solve\_mpcc\_market\_clearing} για την επίλυση.

\subsection{Ροή Εργασιών}
\label{subsec:workflow}

Η κεντρική ροή εργασιών του συστήματος για τον υπολογισμό τιμών αγοράς περιλαμβάνει τα ακόλουθα βήματα. Αρχικά, γίνεται φόρτωση δεδομένων από τη βάση, συμπεριλαμβανομένων των διαθέσιμων μονάδων παραγωγής (με φιλτράρισμα μη διαθεσιμοτήτων), της ζήτησης για την επιλεγμένη ημερομηνία, και των προβλέψεων παραγωγής ΑΠΕ. Στη συνέχεια, επιλύεται το πρόβλημα Unit Commitment για τον προσδιορισμό του βέλτιστου προγράμματος λειτουργίας. Ακολουθεί η μετατροπή της λύσης του UC σε εντολές αγοράς σύμφωνα με την επιλεγμένη στρατηγική. Εκτελείται η εκκαθάριση της αγοράς μέσω του αλγορίθμου MPCC, που καθορίζει τις τιμές ισορροπίας. Τέλος, τα αποτελέσματα αποθηκεύονται στη βάση δεδομένων.

Η κύρια συνάρτηση για την εκτέλεση αυτής της ροής είναι η \texttt{generate\_energy\_prices}:

\begin{samepage}
\begin{minted}[fontsize=\small,linenos,frame=lines]{julia}
# Παράδειγμα χρήσης για μονοζωνική εκκαθάριση
prices = generate_energy_prices("GR", Date(2024, 6, 15);
    order_method = :uc_based,
    markup_factor = 1.1,
    save_to_db = true
)

# Πολυζωνική εκκαθάριση με περιορισμούς ATC
result = run_multi_zone_market_clearing(Date(2024, 6, 15);
    zones = ["GR", "BG", "RO"],
    order_method = :alternative,
    save_to_db = true
)
\end{minted}
\end{samepage}

\subsection{Διαχείριση Επιλυτών}
\label{subsec:solver-management}

Το σύστημα υποστηρίζει πολλαπλούς επιλυτές μέσω ενός μηχανισμού αυτόματης επιλογής και caching. Η συνάρτηση \texttt{select\_solver} ελέγχει τη διαθεσιμότητα των εγκατεστημένων επιλυτών (HiGHS, Gurobi, CPLEX) και επιλέγει τον κατάλληλο βάσει προτίμησης ή διαθεσιμότητας.

Για τον επιλυτή Gurobi, ο οποίος απαιτεί άδεια χρήσης με σημαντικό κόστος αρχικοποίησης, υλοποιήθηκε ένας μηχανισμός caching του περιβάλλοντος (environment) που αποφεύγει την επανάληψη της διαδικασίας αυθεντικοποίησης σε διαδοχικές κλήσεις.

\begin{samepage}
\begin{minted}[fontsize=\small,linenos,frame=lines]{julia}
# Αυτόματη επιλογή επιλυτή
optimizer, solver_name = select_solver("auto")

# Προτίμηση συγκεκριμένου επιλυτή
optimizer, solver_name = select_solver("gurobi")

# Δημιουργία μοντέλου με τον επιλεγμένο επιλυτή
model = Model(optimizer)
\end{minted}
\end{samepage}

% =============================================================================

\chapter{Αρχιτεκτονική Συστήματος και Δεδομένα}
\label{ch:system-architecture}

Στο παρόν κεφάλαιο περιγράφεται η αρχιτεκτονική του συστήματος που αναπτύχθηκε για την προσομοίωση της Προημερήσιας Αγοράς ηλεκτρικής ενέργειας. Η αρχιτεκτονική βασίζεται σε τρεις πυλώνες: την υποδομή δεδομένων που συλλέγει και αποθηκεύει δεδομένα από την πλατφόρμα διαφάνειας του ENTSO-E, τη βάση δεδομένων PostgreSQL που οργανώνει τα δεδομένα σε κατάλληλο σχήμα, και τον πυρήνα υπολογισμών σε Julia/JuMP που υλοποιεί τους αλγορίθμους βελτιστοποίησης.

Η γενική ροή δεδομένων του συστήματος απεικονίζεται στο Σχήμα~\ref{fig:system-architecture}. Τα δεδομένα συλλέγονται από την πλατφόρμα ENTSO-E μέσω ενός αυτοματοποιημένου pipeline σε Python, αποθηκεύονται στη βάση δεδομένων PostgreSQL, επεξεργάζονται από τους αλγορίθμους Julia/JuMP, και τα αποτελέσματα αποθηκεύονται πίσω στη βάση για ανάλυση και οπτικοποίηση.

% TODO: Εισαγωγή σχήματος αρχιτεκτονικής
\begin{figure}[htbp]
\centering
% \includegraphics[width=0.95\textwidth]{figures/system_architecture.pdf}
\fbox{\parbox{0.9\textwidth}{\centering\vspace{3cm}
\texttt{ENTSO-E Transparency Platform} \\
$\downarrow$ \textit{(Python pipeline)} \\
\texttt{PostgreSQL Database} \\
$\downarrow$ \textit{(Julia/JuMP)} \\
\texttt{Euphemia Engine} \\
$\downarrow$ \\
\texttt{simulations.energy\_prices}
\vspace{1cm}}}
\caption{Γενική αρχιτεκτονική του συστήματος προσομοίωσης αγοράς ηλεκτρικής ενέργειας.}
\label{fig:system-architecture}
\end{figure}


% =============================================================================
\section{Πηγές Δεδομένων}
\label{sec:data-sources}

Η κύρια πηγή δεδομένων για το σύστημα είναι η Πλατφόρμα Διαφάνειας του ENTSO-E (ENTSO-E Transparency Platform), η οποία παρέχει δημοσίως διαθέσιμα δεδομένα για τις ευρωπαϊκές αγορές ηλεκτρικής ενέργειας σύμφωνα με τον Κανονισμό (ΕΕ) 543/2013~\cite{entsoe-transparency, eu-regulation-543}.

\subsection{Πλατφόρμα Διαφάνειας ENTSO-E}
\label{subsec:entsoe-platform}

Το ENTSO-E (European Network of Transmission System Operators for Electricity) είναι ο οργανισμός που συντονίζει τους Διαχειριστές Συστημάτων Μεταφοράς (TSOs) της Ευρώπης. Η Πλατφόρμα Διαφάνειας που διατηρεί παρέχει πρόσβαση σε δεδομένα που καλύπτουν τη λειτουργία των ευρωπαϊκών αγορών ηλεκτρικής ενέργειας, συμπεριλαμβανομένων δεδομένων παραγωγής, κατανάλωσης, τιμών, διασυνοριακών ροών και διαθεσιμότητας μονάδων.

Η πλατφόρμα προσφέρει δύο κύριους τρόπους πρόσβασης. Ο πρώτος είναι το RESTful API, το οποίο επιτρέπει προγραμματιστική πρόσβαση στα δεδομένα μέσω HTTP αιτημάτων. Το API χρησιμοποιεί μορφή XML για τις απαντήσεις και απαιτεί κλειδί αυθεντικοποίησης (API key) που παρέχεται δωρεάν μετά από εγγραφή. Ο δεύτερος τρόπος είναι η SFTP πρόσβαση, η οποία παρέχει πρόσβαση σε προ-επεξεργασμένα αρχεία CSV ανά τύπο δεδομένων. Αυτή η μέθοδος είναι αποτελεσματικότερη για μαζική λήψη δεδομένων και χρησιμοποιείται στην παρούσα υλοποίηση.

\subsection{Τύποι Δεδομένων}
\label{subsec:data-types}

Για την προσομοίωση της Προημερήσιας Αγοράς, το σύστημα συλλέγει και επεξεργάζεται τους ακόλουθους τύπους δεδομένων:

\paragraph{Μονάδες Παραγωγής (Production and Generation Units)}
Ο πίνακας \texttt{production\_and\_generation\_units} περιέχει πληροφορίες για τις μονάδες παραγωγής ηλεκτρικής ενέργειας σε κάθε ζώνη προσφοράς. Τα δεδομένα περιλαμβάνουν τον κωδικό και το όνομα της μονάδας, τη ζώνη προσφοράς (bidding zone), τον τύπο καυσίμου (π.χ. φυσικό αέριο, λιγνίτης, υδροηλεκτρικά, αιολικά), την εγκατεστημένη ισχύ σε MW, και την κατάσταση λειτουργίας (commissioned/decommissioned). Αυτά τα δεδομένα χρησιμοποιούνται για τη δημιουργία του συνόλου των διαθέσιμων μονάδων στο πρόβλημα Unit Commitment.

\paragraph{Πραγματικό Φορτίο (Actual Total Load)}
Ο πίνακας \texttt{actual\_total\_load} περιέχει ιστορικά δεδομένα κατανάλωσης ηλεκτρικής ενέργειας ανά ζώνη προσφοράς με ωριαία ή τεταρτοωριαία ανάλυση. Η ζήτηση αποτελεί βασική παράμετρο εισόδου τόσο στο πρόβλημα Unit Commitment όσο και στην εκκαθάριση της αγοράς.

\paragraph{Προβλέψεις Παραγωγής ΑΠΕ (Generation Forecasts for Wind and Solar)}
Ο πίνακας \texttt{generation\_forecasts\_for\_wind\_and\_solar} περιέχει προβλέψεις παραγωγής από αιολικά και φωτοβολταϊκά συστήματα. Δεδομένου ότι η παραγωγή από ΑΠΕ θεωρείται μη κατανεμόμενη (non-dispatchable), οι προβλέψεις αφαιρούνται από τη συνολική ζήτηση για τον υπολογισμό της καθαρής ζήτησης (net demand) που πρέπει να καλυφθεί από τις θερμικές μονάδες.

\paragraph{Διαθέσιμη Ικανότητα Μεταφοράς (Offered Transfer Capacities)}
Ο πίνακας \texttt{offered\_transfer\_capacities\_implicit} περιέχει δεδομένα για τη διαθέσιμη ικανότητα μεταφοράς (ATC) μεταξύ γειτονικών ζωνών προσφοράς. Τα δεδομένα περιλαμβάνουν τη ζώνη προέλευσης και προορισμού, τη χρονική περίοδο, και τη διαθέσιμη ικανότητα σε MW. Αυτά τα δεδομένα είναι απαραίτητα για την πολυζωνική εκκαθάριση της αγοράς με περιορισμούς δικτύου.

\paragraph{Μη Διαθεσιμότητα Μονάδων (Unavailability of Production Units)}
Ο πίνακας \texttt{unavailability\_of\_production\_and\_generation\_units} περιέχει πληροφορίες για προγραμματισμένες ή απρόβλεπτες διακοπές λειτουργίας μονάδων παραγωγής. Κάθε εγγραφή περιλαμβάνει τον κωδικό της μονάδας, την περίοδο μη διαθεσιμότητας, τον τύπο (προγραμματισμένη ή έκτακτη), την κατάσταση (ενεργή, ακυρωμένη), και την εναπομένουσα διαθέσιμη ισχύ. Αυτά τα δεδομένα χρησιμοποιούνται για τον αποκλεισμό ή τη μείωση της διαθέσιμης ισχύος μονάδων κατά την επίλυση του Unit Commitment.

\paragraph{Πραγματική Παραγωγή ανά Μονάδα (Actual Generation Output)}
Ο πίνακας \texttt{actual\_generation\_output\_per\_generation\_unit} περιέχει ιστορικά δεδομένα πραγματικής παραγωγής ανά μονάδα, με χρονική ανάλυση 15 λεπτών ή 1 ώρας. Κάθε εγγραφή περιλαμβάνει τον κωδικό της μονάδας, τη χρονοσφραγίδα (UTC), τον κωδικό ανάλυσης (\texttt{PT15M}, \texttt{PT60M}) και τη μετρούμενη παραγωγή σε MW. Τα δεδομένα αυτά αξιοποιούνται για τη στατιστική εξαγωγή τεχνικών παραμέτρων (ρυθμοί μεταβολής, ελάχιστη ισχύς, χρόνοι λειτουργίας/αδράνειας) μέσω ανάλυσης ιστορικού παραθύρου 3 μηνών.

\paragraph{Τιμές Ενέργειας (Energy Prices)}
Ο πίνακας \texttt{energy\_prices} περιέχει τις πραγματικές τιμές εκκαθάρισης της Προημερήσιας Αγοράς ανά ζώνη προσφοράς. Αυτά τα δεδομένα χρησιμοποιούνται για την επικύρωση (validation) των αποτελεσμάτων του συστήματος προσομοίωσης.

\subsection{Διαθεσιμότητα και Ποιότητα Δεδομένων}
\label{subsec:data-quality}

Η ποιότητα και η διαθεσιμότητα των δεδομένων του ENTSO-E ποικίλλει ανάλογα με τη ζώνη προσφοράς και τον τύπο δεδομένων. Ορισμένες ζώνες (π.χ. Γερμανία, Γαλλία, Σκανδιναβικές χώρες) διαθέτουν πλήρη και υψηλής ποιότητας δεδομένα, ενώ άλλες μπορεί να έχουν ελλείψεις ή ασυνέπειες. Για την Ελλάδα (ζώνη GR), τα δεδομένα είναι γενικά διαθέσιμα από το 2015 και μετά με καλή ποιότητα.

Οι κύριες προκλήσεις που αντιμετωπίζονται στην επεξεργασία των δεδομένων περιλαμβάνουν διαφορετικές χρονικές αναλύσεις (ωριαία, ημίωρη, τεταρτοωριαία) που απαιτούν ομογενοποίηση, ελλιπείς εγγραφές που απαιτούν συμπλήρωση ή αποκλεισμό, διαφορετικές ζώνες ώρας και μετάβαση θερινής/χειμερινής ώρας, και αλλαγές στη δομή των δεδομένων κατά τη διάρκεια του χρόνου.


% =============================================================================
\section{Αγωγός Δεδομένων (Data Pipeline)}
\label{sec:data-pipeline}

Για τη συστηματική συλλογή και ενημέρωση των δεδομένων από την πλατφόρμα ENTSO-E, αναπτύχθηκε ένας αυτοματοποιημένος αγωγός δεδομένων (data pipeline) σε Python. Ο αγωγός υλοποιεί ένα σύστημα σταδιακών ενημερώσεων (incremental updates) που επιτρέπει την αποτελεσματική ενημέρωση της βάσης δεδομένων χωρίς την ανάγκη πλήρους επαναφόρτωσης~\cite{kleppmann2017}.

\subsection{Αρχιτεκτονική Pipeline}
\label{subsec:pipeline-architecture}

Ο αγωγός δεδομένων αποτελείται από τα ακόλουθα βασικά στοιχεία:

\paragraph{Κεντρική Βιβλιοθήκη (helper\_functions.py)}
Περιέχει τις βασικές λειτουργίες που χρησιμοποιούνται από όλα τα scripts ενημέρωσης, συμπεριλαμβανομένων συναρτήσεων αυθεντικοποίησης με την πλατφόρμα ENTSO-E, ανακάλυψης και φιλτραρίσματος αρχείων, λήψης και επεξεργασίας CSV, μαζικής φόρτωσης στη βάση δεδομένων, και διαχείρισης αρχείων καταγραφής.

\paragraph{Scripts Ενημέρωσης (update\_*.py)}
Για κάθε τύπο δεδομένων υπάρχει ένα ξεχωριστό script ενημέρωσης που καθορίζει τις παραμέτρους ειδικά για αυτόν τον τύπο (όνομα πίνακα, φάκελος ENTSO-E, διαχωριστικό αρχείου, κωδικοποίηση κ.λπ.). Αυτή η αρχιτεκτονική επιτρέπει την εύκολη προσθήκη νέων τύπων δεδομένων και την ανεξάρτητη εκτέλεση κάθε ενημέρωσης.

\paragraph{Εργαλείο Δημιουργίας Scripts (generate\_incremental\_scripts.py)}
Αυτοματοποιεί τη δημιουργία των scripts ενημέρωσης με βάση πρότυπα, διασφαλίζοντας συνέπεια στη δομή και τη λειτουργικότητα.

\subsection{Σύστημα Σταδιακών Ενημερώσεων}
\label{subsec:incremental-updates}

Η κεντρική καινοτομία του αγωγού δεδομένων είναι το σύστημα σταδιακών ενημερώσεων που αποφεύγει την επαναφόρτωση ολόκληρου του συνόλου δεδομένων σε κάθε εκτέλεση. Η ροή εργασιών περιγράφεται ως εξής:

Αρχικά, το σύστημα εκτελεί αυθεντικοποίηση με την πλατφόρμα ENTSO-E χρησιμοποιώντας διαπιστευτήρια που αποθηκεύονται σε μεταβλητές περιβάλλοντος. Στη συνέχεια, συνδέεται μέσω SFTP και ανακαλύπτει τα διαθέσιμα αρχεία στον αντίστοιχο φάκελο, καταγράφοντας το όνομα και την ημερομηνία τελευταίας τροποποίησης κάθε αρχείου. Το σύστημα συγκρίνει τη λίστα με τον πίνακα \texttt{file\_update\_log} της βάσης δεδομένων και εντοπίζει τα νέα ή τροποποιημένα αρχεία. Τα αρχεία επεξεργάζονται σε παρτίδες (batches) καθορισμένου μεγέθους (προεπιλογή: 50 αρχεία) για αποτελεσματική χρήση μνήμης. Κάθε αρχείο CSV αναλύεται, καθαρίζεται (data cleaning) και μετασχηματίζεται στη μορφή του πίνακα προορισμού. Τέλος, τα δεδομένα φορτώνονται στη βάση χρησιμοποιώντας τη μέθοδο PostgreSQL COPY για μέγιστη απόδοση, και ο πίνακας \texttt{file\_update\_log} ενημερώνεται με τα στοιχεία των επεξεργασμένων αρχείων.

\subsection{Βελτιστοποίηση Απόδοσης}
\label{subsec:pipeline-performance}

Για τη βελτιστοποίηση της απόδοσης του αγωγού δεδομένων, υλοποιήθηκαν διάφορες τεχνικές. Η χρήση της μεθόδου PostgreSQL COPY αντί για πολλαπλές εντολές INSERT προσφέρει σημαντικά υψηλότερη απόδοση για μαζική φόρτωση δεδομένων. Το ρυθμιζόμενο μέγεθος παρτίδας (batch size) επιτρέπει την εξισορρόπηση μεταξύ χρήσης μνήμης και απόδοσης. Η επεξεργασία σε παρτίδες αντί της φόρτωσης ολόκληρου του αρχείου στη μνήμη επιτρέπει τη διαχείριση πολύ μεγάλων συνόλων δεδομένων. Ο μηχανισμός παρακολούθησης αρχείων μέσω του \texttt{file\_update\_log} αποφεύγει την περιττή επαναεπεξεργασία αρχείων.

\subsection{Χρονοπρογραμματισμός}
\label{subsec:scheduling}

Η αυτοματοποιημένη εκτέλεση του αγωγού δεδομένων επιτυγχάνεται μέσω GitHub Actions, με τις εργασίες να εκτελούνται σε καθημερινή βάση. Η ροή εργασιών περιλαμβάνει την εκτέλεση του \texttt{generate\_incremental\_scripts.py} για τη δημιουργία των τελευταίων scripts, ακολουθούμενη από την εκτέλεση κάθε script ενημέρωσης για τους διάφορους τύπους δεδομένων.


% =============================================================================
\section{Σχήμα Βάσης Δεδομένων}
\label{sec:database-schema}

Η βάση δεδομένων PostgreSQL οργανώνεται σε δύο κύρια σχήματα (schemas): το σχήμα \texttt{entsoe} που περιέχει τα ακατέργαστα δεδομένα από την πλατφόρμα ENTSO-E, και το σχήμα \texttt{simulations} που περιέχει τα αποτελέσματα των προσομοιώσεων.

\subsection{Σχήμα ENTSO-E}
\label{subsec:entsoe-schema}

Το σχήμα \texttt{entsoe} περιέχει τους πίνακες που αποθηκεύουν τα δεδομένα από την πλατφόρμα διαφάνειας. Οι κύριοι πίνακες παρουσιάζονται στον Πίνακα~\ref{tab:entsoe-tables}.

\begin{table}[htbp]
\centering
\caption{Κύριοι πίνακες του σχήματος \texttt{entsoe}}
\label{tab:entsoe-tables}
\begin{tabular}{@{}p{5.5cm}p{8cm}@{}}
\toprule
\textbf{Πίνακας} & \textbf{Περιγραφή} \\
\midrule
\texttt{production\_and\_generation\_units} & Στοιχεία μονάδων παραγωγής \\
\texttt{actual\_total\_load} & Ιστορικά δεδομένα κατανάλωσης \\
\texttt{actual\_generation\_output\_per\_generation\_unit} & Ιστορική παραγωγή ανά μονάδα \\
\texttt{generation\_forecasts\_for\_wind\_and\_solar} & Προβλέψεις παραγωγής ΑΠΕ \\
\texttt{offered\_transfer\_capacities\_implicit} & Διαθέσιμη ικανότητα μεταφοράς \\
\texttt{unavailability\_of\_production\_and\_generation\_units} & Μη διαθεσιμότητα μονάδων \\
\texttt{energy\_prices} & Τιμές εκκαθάρισης αγοράς \\
\texttt{physical\_flows} & Φυσικές ροές διασυνδέσεων \\
\texttt{file\_update\_log} & Καταγραφή επεξεργασμένων αρχείων \\
\bottomrule
\end{tabular}
\end{table}

\paragraph{Πίνακας production\_and\_generation\_units}
Ο πίνακας αυτός περιέχει τα στοιχεία των μονάδων παραγωγής με τα ακόλουθα βασικά πεδία: \texttt{generation\_unit\_code} (μοναδικός κωδικός μονάδας), \texttt{generation\_unit\_name} (όνομα μονάδας), \texttt{area\_map\_code} (κωδικός ζώνης προσφοράς, π.χ. ``GR''), \texttt{production\_unit\_type} (τύπος καυσίμου), \texttt{generation\_unit\_installed\_capacity\_mw} (εγκατεστημένη ισχύς), \texttt{production\_unit\_status} (κατάσταση λειτουργίας), \texttt{valid\_from} και \texttt{valid\_to} (περίοδος ισχύος των στοιχείων).

Η ερώτηση για την ανάκτηση διαθέσιμων μονάδων για μια συγκεκριμένη ζώνη και ημερομηνία φιλτράρει τις μονάδες με βάση την κατάσταση \texttt{COMMISSIONED}, τον τύπο περιοχής \texttt{BZN} ή \texttt{BZN/CTA}, και την περίοδο ισχύος.

\paragraph{Πίνακας actual\_total\_load}
Ο πίνακας αυτός περιέχει χρονοσειρές κατανάλωσης με τα πεδία: \texttt{datetime\_utc} (χρονοσφραγίδα σε UTC), \texttt{area\_map\_code} (ζώνη προσφοράς), \texttt{total\_load\_mw} (κατανάλωση σε MW), και \texttt{resolution} (χρονική ανάλυση).

\paragraph{Πίνακας actual\_generation\_output\_per\_generation\_unit}
Ο πίνακας αυτός περιέχει ιστορικά δεδομένα πραγματικής παραγωγής ανά μονάδα με τα πεδία: \texttt{generation\_unit\_code} (κωδικός μονάδας), \texttt{datetime\_utc} (χρονοσφραγίδα σε UTC), \texttt{actual\_generation\_output\_mw} (μετρούμενη παραγωγή σε MW), και \texttt{resolution} (χρονική ανάλυση, συνήθως 15 λεπτά ή 1 ώρα). Τα δεδομένα αυτά χρησιμοποιούνται για την εξαγωγή τεχνικών παραμέτρων (ρυθμοί μεταβολής, ελάχιστη ισχύς, χρόνοι λειτουργίας/αδράνειας) μέσω στατιστικής ανάλυσης, όπως περιγράφεται στην Ενότητα~\ref{subsec:parameter-inference}.

\paragraph{Πίνακας unavailability\_of\_production\_and\_generation\_units}
Ο πίνακας αυτός περιέχει πληροφορίες διακοπών με τα πεδία: \texttt{asset\_code} (αντιστοιχεί στο \texttt{generation\_unit\_code}), \texttt{start\_outage\_utc} και \texttt{end\_outage\_utc} (περίοδος διακοπής), \texttt{type} (``Planned'' ή ``Forced''), \texttt{status} (``Active'', ``Cancelled'', ``Withdrawn''), και \texttt{available\_capacity\_mw} (εναπομένουσα διαθέσιμη ισχύς).

\subsection{Σχήμα Προσομοιώσεων}
\label{subsec:simulations-schema}

Το σχήμα \texttt{simulations} περιέχει τους πίνακες για την αποθήκευση των αποτελεσμάτων των προσομοιώσεων. Οι κύριοι πίνακες παρουσιάζονται στον Πίνακα~\ref{tab:simulations-tables}.

\begin{table}[htbp]
\centering
\caption{Κύριοι πίνακες του σχήματος \texttt{simulations}}
\label{tab:simulations-tables}
\begin{tabular}{@{}p{5cm}p{8cm}@{}}
\toprule
\textbf{Πίνακας} & \textbf{Περιγραφή} \\
\midrule
\texttt{energy\_prices} & Υπολογισμένες τιμές εκκαθάρισης \\
\texttt{optimization\_runs} & Μεταδεδομένα εκτελέσεων βελτιστοποίησης \\
\texttt{transmission\_flows} & Διασυνοριακές ροές ισχύος \\
\texttt{generator\_inferred\_parameters} & Εξαχθείσες τεχνικές παράμετροι μονάδων \\
\texttt{uc\_results} & Αποτελέσματα Unit Commitment (σύνοψη) \\
\texttt{uc\_generation} & Παραγωγή/δέσμευση ανά μονάδα και περίοδο \\
\texttt{uc\_net\_demand} & Καθαρή ζήτηση ανά περίοδο \\
\bottomrule
\end{tabular}
\end{table}

\paragraph{Πίνακας energy\_prices}
Ο πίνακας αυτός αποθηκεύει τις υπολογισμένες τιμές εκκαθάρισης με τα πεδία: \texttt{bidding\_zone} (ζώνη προσφοράς), \texttt{date} (ημερομηνία), \texttt{hour} (ώρα 1-24), \texttt{price} (τιμή σε €/MWh), \texttt{created\_at} (χρονοσφραγίδα δημιουργίας), και \texttt{run\_id} (αναφορά στην εκτέλεση βελτιστοποίησης).

\paragraph{Πίνακας transmission\_flows}
Ο πίνακας αυτός αποθηκεύει τις υπολογισμένες διασυνοριακές ροές με τα πεδία: \texttt{source\_zone} και \texttt{target\_zone} (ζώνες προέλευσης και προορισμού), \texttt{date} (ημερομηνία), \texttt{hour} (ώρα), \texttt{flow\_mw} (ροή σε MW), και \texttt{run\_id}.

\paragraph{Πίνακας generator\_inferred\_parameters}
Ο πίνακας αυτός λειτουργεί ως κρυφή μνήμη (cache) για τις τεχνικές παραμέτρους μονάδων παραγωγής που εξάγονται από ιστορικά δεδομένα. Τα βασικά πεδία περιλαμβάνουν: \texttt{generator\_code} (κωδικός μονάδας), \texttt{ramp\_up} και \texttt{ramp\_down} (ρυθμοί μεταβολής ως κλάσμα της $P^{\max}$ ανά ώρα), \texttt{p\_min} (ελάχιστη ισχύς ως κλάσμα της $P^{\max}$), \texttt{min\_uptime} και \texttt{min\_downtime} (ελάχιστοι χρόνοι λειτουργίας/αδράνειας σε ώρες), και \texttt{inference\_date} (ημερομηνία εξαγωγής για έλεγχο λήξης). Η χρήση του πίνακα αυτού μειώνει τον χρόνο ανάκτησης παραμέτρων από περίπου 17 λεπτά σε 2 δευτερόλεπτα.

\paragraph{Πίνακες uc\_results, uc\_generation, uc\_net\_demand}
Οι τρεις αυτοί πίνακες αποτελούν τον μηχανισμό κρυφής μνήμης (caching) για τα αποτελέσματα του επιλυτή Unit Commitment. Ο πίνακας \texttt{uc\_results} αποθηκεύει τη σύνοψη κάθε εκτέλεσης με κλειδί \texttt{(bidding\_zone, market\_date, code\_version)} και περιλαμβάνει τα πεδία: \texttt{status} (κατάσταση τερματισμού), \texttt{solver} (επιλυτής), \texttt{resolution\_minutes} (χρονική ανάλυση), \texttt{total\_cost}, \texttt{production\_cost}, \texttt{startup\_cost}, \texttt{noload\_cost} (συνιστώσες κόστους), και αριθμούς εκκινήσεων (\texttt{hot/warm/cold\_startups}). Ο πίνακας \texttt{uc\_generation} αποθηκεύει τα αναλυτικά αποτελέσματα παραγωγής, δέσμευσης και εκκίνησης ανά μονάδα και περίοδο, ενώ ο πίνακας \texttt{uc\_net\_demand} αποθηκεύει την καθαρή ζήτηση και την παραγωγή ΑΠΕ ανά περίοδο. Η αποθήκευση ανέρχεται σε περίπου 195 KB ανά ζώνη/ημέρα.


% =============================================================================
\section{Οικοσύστημα Julia/JuMP}
\label{sec:julia-jump}

Για την υλοποίηση των αλγορίθμων βελτιστοποίησης επιλέχθηκε η γλώσσα προγραμματισμού Julia σε συνδυασμό με το πλαίσιο μοντελοποίησης JuMP~\cite{bezanson2017, jump2017}. Αυτή η επιλογή βασίζεται σε τεχνικά πλεονεκτήματα που καθιστούν το συγκεκριμένο οικοσύστημα ιδανικό για επιστημονικούς υπολογισμούς και προβλήματα βελτιστοποίησης.

\subsection{Γλώσσα Προγραμματισμού Julia}
\label{subsec:julia-language}

Η Julia είναι μια σύγχρονη γλώσσα προγραμματισμού υψηλού επιπέδου σχεδιασμένη για επιστημονικούς υπολογισμούς. Τα κύρια χαρακτηριστικά της που την καθιστούν κατάλληλη για την παρούσα εφαρμογή είναι τα ακόλουθα.

Η Julia επιτυγχάνει υψηλή απόδοση μέσω της μεταγλώττισης just-in-time (JIT) με χρήση του LLVM compiler framework. Ο κώδικας Julia μεταγλωττίζεται σε εγγενή κώδικα μηχανής, επιτυγχάνοντας ταχύτητες συγκρίσιμες με C και Fortran, χωρίς την ανάγκη για χειροκίνητη βελτιστοποίηση ή κλήση εξωτερικών βιβλιοθηκών.

Παρά την υψηλή απόδοση, η σύνταξη της Julia είναι εκφραστική και φιλική στον χρήστη, παρόμοια με Python ή MATLAB. Αυτό επιτρέπει την ταχεία ανάπτυξη πρωτοτύπων και την εύκολη ανάγνωση του κώδικα.

Η Julia διαθέτει ένα εξελιγμένο σύστημα τύπων με πολλαπλή αποστολή (multiple dispatch) που επιτρέπει τη δημιουργία γενικευμένου κώδικα που είναι ταυτόχρονα αποτελεσματικός και επαναχρησιμοποιήσιμος.

Η Julia παρέχει ενσωματωμένη υποστήριξη για παραλληλισμό και κατανεμημένους υπολογισμούς, διευκολύνοντας την κλιμάκωση σε πολλαπλούς πυρήνες ή μηχανές.

\subsection{Πλαίσιο Μοντελοποίησης JuMP}
\label{subsec:jump-framework}

Το JuMP (Julia for Mathematical Programming) είναι ένα domain-specific language (DSL) ενσωματωμένο στη Julia για τη διατύπωση και επίλυση προβλημάτων μαθηματικού προγραμματισμού. Αποτελεί ένα από τα πιο δημοφιλή και ώριμα πλαίσια μοντελοποίησης βελτιστοποίησης στον ακαδημαϊκό και βιομηχανικό χώρο~\cite{jump2017}.

Τα πλεονεκτήματα του JuMP περιλαμβάνουν μια αλγεβρική μοντελοποίηση που επιτρέπει τη διατύπωση προβλημάτων βελτιστοποίησης με τρόπο που μοιάζει με τη μαθηματική τους διατύπωση. Η σύνταξη για μεταβλητές, περιορισμούς και αντικειμενικές συναρτήσεις είναι διαισθητική και εκφραστική. Το JuMP υποστηρίζει πολλαπλούς επιλυτές (solvers) μέσω μιας ενιαίας διεπαφής, επιτρέποντας την εύκολη εναλλαγή μεταξύ επιλυτών χωρίς αλλαγές στον κώδικα μοντελοποίησης. Η απόδοση του JuMP είναι συγκρίσιμη ή καλύτερη από εμπορικά εργαλεία μοντελοποίησης όπως το AMPL και το GAMS. Τέλος, το JuMP είναι εντελώς ελεύθερο και ανοιχτού κώδικα, διατίθεται υπό την άδεια MIT.

Ένα τυπικό παράδειγμα χρήσης του JuMP για τη δημιουργία ενός μοντέλου βελτιστοποίησης φαίνεται στο ακόλουθο απόσπασμα κώδικα:

\begin{samepage}
\begin{minted}[fontsize=\small,linenos,frame=lines]{julia}
using JuMP, HiGHS

# Δημιουργία μοντέλου με τον επιλυτή HiGHS
model = Model(HiGHS.Optimizer)

# Ορισμός μεταβλητών
@variable(model, x >= 0)
@variable(model, y >= 0)

# Ορισμός περιορισμών
@constraint(model, x + 2y <= 10)
@constraint(model, 3x + y <= 15)

# Ορισμός αντικειμενικής συνάρτησης
@objective(model, Max, 2x + 3y)

# Επίλυση
optimize!(model)
\end{minted}
\end{samepage}

\subsection{Επιλυτής HiGHS}
\label{subsec:highs-solver}

Ως κύριος επιλυτής για τα προβλήματα βελτιστοποίησης χρησιμοποιείται ο HiGHS, ένας σύγχρονος επιλυτής ανοιχτού κώδικα για προβλήματα γραμμικού προγραμματισμού (LP), μικτού ακέραιου προγραμματισμού (MIP) και τετραγωνικού προγραμματισμού (QP)~\cite{highs2024}.

Ο HiGHS αναπτύσσεται στο Πανεπιστήμιο του Εδιμβούργου και αποτελεί έναν από τους ταχύτερους ανοιχτού κώδικα επιλυτές. Περιλαμβάνει υλοποιήσεις του αλγορίθμου dual revised simplex με εκτεταμένη παραλληλοποίηση, μεθόδων εσωτερικού σημείου (interior point methods), και αλγορίθμων branch-and-bound για MIP με cutting planes και heuristics.

Η διεπαφή με τη Julia παρέχεται μέσω του πακέτου HiGHS.jl, το οποίο ενσωματώνεται απρόσκοπτα με το JuMP. Η επιλογή του HiGHS ως προεπιλεγμένου επιλυτή βασίζεται στην άδεια ανοιχτού κώδικα (MIT) που επιτρέπει ελεύθερη χρήση, στην υψηλή απόδοση συγκρίσιμη με εμπορικούς επιλυτές, στην ενεργή ανάπτυξη και υποστήριξη, και στη συμβατότητα με όλα τα απαιτούμενα είδη προβλημάτων (LP, MILP).

Το σύστημα υποστηρίζει επίσης εμπορικούς επιλυτές όπως ο Gurobi και ο CPLEX για περιπτώσεις που απαιτείται ακόμη υψηλότερη απόδοση, με την εναλλαγή να γίνεται μέσω μιας απλής παραμέτρου.

\subsection{Βασικά Πακέτα Julia}
\label{subsec:julia-packages}

Εκτός από το JuMP και τον HiGHS, το σύστημα χρησιμοποιεί μια σειρά από πακέτα Julia για διάφορες λειτουργίες. Το πακέτο DataFrames.jl χρησιμοποιείται για τη διαχείριση και επεξεργασία δεδομένων σε μορφή πινάκων, με λειτουργικότητα παρόμοια με τα pandas της Python. Το πακέτο LibPQ.jl παρέχει σύνδεση με τη βάση δεδομένων PostgreSQL μέσω της βιβλιοθήκης libpq, υποστηρίζοντας παραμετροποιημένα ερωτήματα και διαχείριση συνδέσεων. Το πακέτο CSV.jl χρησιμοποιείται για την ανάγνωση και εγγραφή αρχείων CSV με υψηλή απόδοση. Το πακέτο Dates.jl, που αποτελεί μέρος της standard library, χρησιμοποιείται για τη διαχείριση ημερομηνιών και χρόνων. Το πακέτο Statistics.jl, επίσης μέρος της standard library, χρησιμοποιείται για στατιστικούς υπολογισμούς---κυρίως τη συνάρτηση \texttt{quantile} για τον υπολογισμό εκατοστημορίων κατά την εξαγωγή τεχνικών παραμέτρων μονάδων από ιστορικά δεδομένα. Τέλος, το πακέτο DotEnv.jl χρησιμοποιείται για τη φόρτωση μεταβλητών περιβάλλοντος από αρχεία \texttt{.env} για τη διαχείριση διαπιστευτηρίων.


% =============================================================================
\section{Αρχιτεκτονική Εφαρμογής}
\label{sec:application-architecture}

Η εφαρμογή Julia οργανώνεται ως ένα ενιαίο module ονόματι \texttt{Euphemia} που περιέχει όλη τη λειτουργικότητα του συστήματος. Η δομή του module ακολουθεί τις βέλτιστες πρακτικές της Julia για οργάνωση κώδικα και επαναχρησιμοποίηση~\cite{julia-docs}.

\subsection{Δομή Module}
\label{subsec:module-structure}

Το module \texttt{Euphemia} αποτελείται από τα ακόλουθα υπο-modules και αρχεία, όπως φαίνεται στο Σχήμα~\ref{fig:module-structure}:

\begin{figure}[htbp]
\centering
\begin{verbatim}
Euphemia.jl (Main Module)
├── dbutils.jl                 # Database utilities
├── Generators.jl              # Generator data model
├── Loads.jl                   # Load data model
├── Renewables.jl              # RES forecasts
├── FuelTypeParameters.jl      # Fuel-specific parameters
├── MarketOrders.jl            # Order types
├── UnitCommitment.jl          # UC solver
├── BiddingStrategy.jl         # UC → Market orders
├── Network.jl                 # ATC constraints
├── MPCC.jl                    # Market clearing
└── AlternativeOrderBook.jl    # Fast order generation
\end{verbatim}
\caption{Δομή του module \texttt{Euphemia}.}
\label{fig:module-structure}
\end{figure}

\paragraph{Κεντρικό Module (Euphemia.jl)}
Αποτελεί το σημείο εισόδου της εφαρμογής. Διαχειρίζεται τις εξαρτήσεις, συντονίζει τα υπο-modules, και εξάγει τη δημόσια διεπαφή (public API). Περιλαμβάνει επίσης τη λογική επιλογής επιλυτή και την κεντρική ροή εργασιών.

\paragraph{Βοηθητικά Βάσης Δεδομένων (dbutils.jl)}
Περιέχει συναρτήσεις για τη σύνδεση με τη βάση δεδομένων PostgreSQL και την εκτέλεση ερωτημάτων. Υλοποιεί το pattern \texttt{withdb} για ασφαλή διαχείριση συνδέσεων με αυτόματο κλείσιμο.

\paragraph{Μοντέλα Δεδομένων (Generators.jl, Loads.jl, Renewables.jl)}
Ορίζουν τις δομές δεδομένων για τις οντότητες του συστήματος και παρέχουν συναρτήσεις για την ανάκτησή τους από τη βάση δεδομένων. Η δομή \texttt{Generator} περιλαμβάνει πεδία για τον κωδικό, το όνομα, τη ζώνη προσφοράς, τον τύπο καυσίμου, τα όρια ισχύος, το οριακό κόστος, και προαιρετικούς τεχνικούς περιορισμούς (ρυθμοί μεταβολής, ελάχιστοι χρόνοι λειτουργίας/αδράνειας) που εξάγονται αυτόματα από ιστορικά δεδομένα. Η δομή \texttt{InitialConditions} αναπαριστά την αρχική κατάσταση κάθε μονάδας στη χρονική στιγμή $t=0$ (δέσμευση, ισχύς εξόδου, θερμοκρασιακή κατάσταση) και εξάγεται από ιστορικά δεδομένα 72 ωρών μέσω μαζικού ερωτήματος. Το αρχείο \texttt{Generators.jl} περιλαμβάνει τις συναρτήσεις εξαγωγής παραμέτρων (\texttt{infer\_ramp\_rates}, \texttt{infer\_p\_min}, \texttt{infer\_uptime\_downtime}), τη συνάρτηση \texttt{get\_initial\_conditions} για αυτόματο προσδιορισμό αρχικών συνθηκών, και τον μηχανισμό αποθήκευσης σε κρυφή μνήμη βάσης δεδομένων.

\paragraph{Τύποι Εντολών (MarketOrders.jl)}
Ορίζει τις δομές δεδομένων για τους διάφορους τύπους εντολών αγοράς, συμπεριλαμβανομένων των \texttt{SimpleOrder}, \texttt{BlockOrder}, \texttt{LinkedBlockOrder} κ.λπ.

\paragraph{Unit Commitment (UnitCommitment.jl)}
Υλοποιεί τον επιλυτή για το πρόβλημα Unit Commitment χρησιμοποιώντας JuMP/HiGHS. Περιλαμβάνει τη διατύπωση του μοντέλου με τριμερή αντικειμενική συνάρτηση (κόστος παραγωγής, κόστος εκκίνησης ανά θερμοκρασιακή κατάσταση, κόστος αδράνειας), την ενσωμάτωση αρχικών συνθηκών, ρυθμιζόμενες παραμέτρους επιλυτή (MIP gap, χρονικό όριο μέσω MOI), καθώς και μηχανισμό αποθήκευσης αποτελεσμάτων σε κρυφή μνήμη βάσης δεδομένων.

\paragraph{Στρατηγική Προσφορών (BiddingStrategy.jl)}
Υλοποιεί τη μετατροπή των αποτελεσμάτων του Unit Commitment σε εντολές αγοράς σύμφωνα με τις τρεις στρατηγικές που περιγράφηκαν στο Κεφάλαιο~\ref{sec:bidding-strategies}.

\paragraph{Δίκτυο (Network.jl)}
Διαχειρίζεται τους περιορισμούς δικτύου και τη διαθέσιμη ικανότητα μεταφοράς (ATC) μεταξύ ζωνών προσφοράς. Παρέχει λειτουργίες για την ανάκτηση δεδομένων ATC από το ENTSO-E και την προσθήκη περιορισμών στο μοντέλο βελτιστοποίησης.

\paragraph{Εκκαθάριση Αγοράς (MPCC.jl)}
Υλοποιεί τον αλγόριθμο εκκαθάρισης της αγοράς με τη μέθοδο MPCC. Περιλαμβάνει τη δομή \texttt{MPCCOrderBook} για τη διαχείριση των εντολών και τη συνάρτηση \texttt{solve\_mpcc\_market\_clearing} για την επίλυση.

\subsection{Ροή Εργασιών}
\label{subsec:workflow}

Η κεντρική ροή εργασιών του συστήματος για τον υπολογισμό τιμών αγοράς περιλαμβάνει τα ακόλουθα βήματα. Αρχικά, γίνεται φόρτωση δεδομένων από τη βάση, συμπεριλαμβανομένων των διαθέσιμων μονάδων παραγωγής (με φιλτράρισμα μη διαθεσιμοτήτων), της ζήτησης για την επιλεγμένη ημερομηνία, και των προβλέψεων παραγωγής ΑΠΕ. Στη συνέχεια, επιλύεται το πρόβλημα Unit Commitment για τον προσδιορισμό του βέλτιστου προγράμματος λειτουργίας. Ακολουθεί η μετατροπή της λύσης του UC σε εντολές αγοράς σύμφωνα με την επιλεγμένη στρατηγική. Εκτελείται η εκκαθάριση της αγοράς μέσω του αλγορίθμου MPCC, που καθορίζει τις τιμές ισορροπίας. Τέλος, τα αποτελέσματα αποθηκεύονται στη βάση δεδομένων.

Η κύρια συνάρτηση για την εκτέλεση αυτής της ροής είναι η \texttt{generate\_energy\_prices}:

\begin{samepage}
\begin{minted}[fontsize=\small,linenos,frame=lines]{julia}
# Παράδειγμα χρήσης για μονοζωνική εκκαθάριση
prices = generate_energy_prices("GR", Date(2024, 6, 15);
    order_method = :uc_based,
    markup_factor = 1.1,
    save_to_db = true
)

# Πολυζωνική εκκαθάριση με περιορισμούς ATC
result = run_multi_zone_market_clearing(Date(2024, 6, 15);
    zones = ["GR", "BG", "RO"],
    order_method = :alternative,
    save_to_db = true
)
\end{minted}
\end{samepage}

\subsection{Διαχείριση Επιλυτών}
\label{subsec:solver-management}

Το σύστημα υποστηρίζει πολλαπλούς επιλυτές μέσω ενός μηχανισμού αυτόματης επιλογής και caching. Η συνάρτηση \texttt{select\_solver} ελέγχει τη διαθεσιμότητα των εγκατεστημένων επιλυτών (HiGHS, Gurobi, CPLEX) και επιλέγει τον κατάλληλο βάσει προτίμησης ή διαθεσιμότητας.

Για τον επιλυτή Gurobi, ο οποίος απαιτεί άδεια χρήσης με σημαντικό κόστος αρχικοποίησης, υλοποιήθηκε ένας μηχανισμός caching του περιβάλλοντος (environment) που αποφεύγει την επανάληψη της διαδικασίας αυθεντικοποίησης σε διαδοχικές κλήσεις.

\begin{samepage}
\begin{minted}[fontsize=\small,linenos,frame=lines]{julia}
# Αυτόματη επιλογή επιλυτή
optimizer, solver_name = select_solver("auto")

# Προτίμηση συγκεκριμένου επιλυτή
optimizer, solver_name = select_solver("gurobi")

# Δημιουργία μοντέλου με τον επιλεγμένο επιλυτή
model = Model(optimizer)
\end{minted}
\end{samepage}

% =============================================================================

\chapter{Αρχιτεκτονική Συστήματος και Δεδομένα}
\label{ch:system-architecture}

Στο παρόν κεφάλαιο περιγράφεται η αρχιτεκτονική του συστήματος που αναπτύχθηκε για την προσομοίωση της Προημερήσιας Αγοράς ηλεκτρικής ενέργειας. Η αρχιτεκτονική βασίζεται σε τρεις πυλώνες: την υποδομή δεδομένων που συλλέγει και αποθηκεύει δεδομένα από την πλατφόρμα διαφάνειας του ENTSO-E, τη βάση δεδομένων PostgreSQL που οργανώνει τα δεδομένα σε κατάλληλο σχήμα, και τον πυρήνα υπολογισμών σε Julia/JuMP που υλοποιεί τους αλγορίθμους βελτιστοποίησης.

Η γενική ροή δεδομένων του συστήματος απεικονίζεται στο Σχήμα~\ref{fig:system-architecture}. Τα δεδομένα συλλέγονται από την πλατφόρμα ENTSO-E μέσω ενός αυτοματοποιημένου pipeline σε Python, αποθηκεύονται στη βάση δεδομένων PostgreSQL, επεξεργάζονται από τους αλγορίθμους Julia/JuMP, και τα αποτελέσματα αποθηκεύονται πίσω στη βάση για ανάλυση και οπτικοποίηση.

% TODO: Εισαγωγή σχήματος αρχιτεκτονικής
\begin{figure}[htbp]
\centering
% \includegraphics[width=0.95\textwidth]{figures/system_architecture.pdf}
\fbox{\parbox{0.9\textwidth}{\centering\vspace{3cm}
\texttt{ENTSO-E Transparency Platform} \\
$\downarrow$ \textit{(Python pipeline)} \\
\texttt{PostgreSQL Database} \\
$\downarrow$ \textit{(Julia/JuMP)} \\
\texttt{Euphemia Engine} \\
$\downarrow$ \\
\texttt{simulations.energy\_prices}
\vspace{1cm}}}
\caption{Γενική αρχιτεκτονική του συστήματος προσομοίωσης αγοράς ηλεκτρικής ενέργειας.}
\label{fig:system-architecture}
\end{figure}


% =============================================================================
\section{Πηγές Δεδομένων}
\label{sec:data-sources}

Η κύρια πηγή δεδομένων για το σύστημα είναι η Πλατφόρμα Διαφάνειας του ENTSO-E (ENTSO-E Transparency Platform), η οποία παρέχει δημοσίως διαθέσιμα δεδομένα για τις ευρωπαϊκές αγορές ηλεκτρικής ενέργειας σύμφωνα με τον Κανονισμό (ΕΕ) 543/2013~\cite{entsoe-transparency, eu-regulation-543}.

\subsection{Πλατφόρμα Διαφάνειας ENTSO-E}
\label{subsec:entsoe-platform}

Το ENTSO-E (European Network of Transmission System Operators for Electricity) είναι ο οργανισμός που συντονίζει τους Διαχειριστές Συστημάτων Μεταφοράς (TSOs) της Ευρώπης. Η Πλατφόρμα Διαφάνειας που διατηρεί παρέχει πρόσβαση σε δεδομένα που καλύπτουν τη λειτουργία των ευρωπαϊκών αγορών ηλεκτρικής ενέργειας, συμπεριλαμβανομένων δεδομένων παραγωγής, κατανάλωσης, τιμών, διασυνοριακών ροών και διαθεσιμότητας μονάδων.

Η πλατφόρμα προσφέρει δύο κύριους τρόπους πρόσβασης. Ο πρώτος είναι το RESTful API, το οποίο επιτρέπει προγραμματιστική πρόσβαση στα δεδομένα μέσω HTTP αιτημάτων. Το API χρησιμοποιεί μορφή XML για τις απαντήσεις και απαιτεί κλειδί αυθεντικοποίησης (API key) που παρέχεται δωρεάν μετά από εγγραφή. Ο δεύτερος τρόπος είναι η SFTP πρόσβαση, η οποία παρέχει πρόσβαση σε προ-επεξεργασμένα αρχεία CSV ανά τύπο δεδομένων. Αυτή η μέθοδος είναι αποτελεσματικότερη για μαζική λήψη δεδομένων και χρησιμοποιείται στην παρούσα υλοποίηση.

\subsection{Τύποι Δεδομένων}
\label{subsec:data-types}

Για την προσομοίωση της Προημερήσιας Αγοράς, το σύστημα συλλέγει και επεξεργάζεται τους ακόλουθους τύπους δεδομένων:

\paragraph{Μονάδες Παραγωγής (Production and Generation Units)}
Ο πίνακας \texttt{production\_and\_generation\_units} περιέχει πληροφορίες για τις μονάδες παραγωγής ηλεκτρικής ενέργειας σε κάθε ζώνη προσφοράς. Τα δεδομένα περιλαμβάνουν τον κωδικό και το όνομα της μονάδας, τη ζώνη προσφοράς (bidding zone), τον τύπο καυσίμου (π.χ. φυσικό αέριο, λιγνίτης, υδροηλεκτρικά, αιολικά), την εγκατεστημένη ισχύ σε MW, και την κατάσταση λειτουργίας (commissioned/decommissioned). Αυτά τα δεδομένα χρησιμοποιούνται για τη δημιουργία του συνόλου των διαθέσιμων μονάδων στο πρόβλημα Unit Commitment.

\paragraph{Πραγματικό Φορτίο (Actual Total Load)}
Ο πίνακας \texttt{actual\_total\_load} περιέχει ιστορικά δεδομένα κατανάλωσης ηλεκτρικής ενέργειας ανά ζώνη προσφοράς με ωριαία ή τεταρτοωριαία ανάλυση. Η ζήτηση αποτελεί βασική παράμετρο εισόδου τόσο στο πρόβλημα Unit Commitment όσο και στην εκκαθάριση της αγοράς.

\paragraph{Προβλέψεις Παραγωγής ΑΠΕ (Generation Forecasts for Wind and Solar)}
Ο πίνακας \texttt{generation\_forecasts\_for\_wind\_and\_solar} περιέχει προβλέψεις παραγωγής από αιολικά και φωτοβολταϊκά συστήματα. Δεδομένου ότι η παραγωγή από ΑΠΕ θεωρείται μη κατανεμόμενη (non-dispatchable), οι προβλέψεις αφαιρούνται από τη συνολική ζήτηση για τον υπολογισμό της καθαρής ζήτησης (net demand) που πρέπει να καλυφθεί από τις θερμικές μονάδες.

\paragraph{Διαθέσιμη Ικανότητα Μεταφοράς (Offered Transfer Capacities)}
Ο πίνακας \texttt{offered\_transfer\_capacities\_implicit} περιέχει δεδομένα για τη διαθέσιμη ικανότητα μεταφοράς (ATC) μεταξύ γειτονικών ζωνών προσφοράς. Τα δεδομένα περιλαμβάνουν τη ζώνη προέλευσης και προορισμού, τη χρονική περίοδο, και τη διαθέσιμη ικανότητα σε MW. Αυτά τα δεδομένα είναι απαραίτητα για την πολυζωνική εκκαθάριση της αγοράς με περιορισμούς δικτύου.

\paragraph{Μη Διαθεσιμότητα Μονάδων (Unavailability of Production Units)}
Ο πίνακας \texttt{unavailability\_of\_production\_and\_generation\_units} περιέχει πληροφορίες για προγραμματισμένες ή απρόβλεπτες διακοπές λειτουργίας μονάδων παραγωγής. Κάθε εγγραφή περιλαμβάνει τον κωδικό της μονάδας, την περίοδο μη διαθεσιμότητας, τον τύπο (προγραμματισμένη ή έκτακτη), την κατάσταση (ενεργή, ακυρωμένη), και την εναπομένουσα διαθέσιμη ισχύ. Αυτά τα δεδομένα χρησιμοποιούνται για τον αποκλεισμό ή τη μείωση της διαθέσιμης ισχύος μονάδων κατά την επίλυση του Unit Commitment.

\paragraph{Πραγματική Παραγωγή ανά Μονάδα (Actual Generation Output)}
Ο πίνακας \texttt{actual\_generation\_output\_per\_generation\_unit} περιέχει ιστορικά δεδομένα πραγματικής παραγωγής ανά μονάδα, με χρονική ανάλυση 15 λεπτών ή 1 ώρας. Κάθε εγγραφή περιλαμβάνει τον κωδικό της μονάδας, τη χρονοσφραγίδα (UTC), τον κωδικό ανάλυσης (\texttt{PT15M}, \texttt{PT60M}) και τη μετρούμενη παραγωγή σε MW. Τα δεδομένα αυτά αξιοποιούνται για τη στατιστική εξαγωγή τεχνικών παραμέτρων (ρυθμοί μεταβολής, ελάχιστη ισχύς, χρόνοι λειτουργίας/αδράνειας) μέσω ανάλυσης ιστορικού παραθύρου 3 μηνών.

\paragraph{Τιμές Ενέργειας (Energy Prices)}
Ο πίνακας \texttt{energy\_prices} περιέχει τις πραγματικές τιμές εκκαθάρισης της Προημερήσιας Αγοράς ανά ζώνη προσφοράς. Αυτά τα δεδομένα χρησιμοποιούνται για την επικύρωση (validation) των αποτελεσμάτων του συστήματος προσομοίωσης.

\subsection{Διαθεσιμότητα και Ποιότητα Δεδομένων}
\label{subsec:data-quality}

Η ποιότητα και η διαθεσιμότητα των δεδομένων του ENTSO-E ποικίλλει ανάλογα με τη ζώνη προσφοράς και τον τύπο δεδομένων. Ορισμένες ζώνες (π.χ. Γερμανία, Γαλλία, Σκανδιναβικές χώρες) διαθέτουν πλήρη και υψηλής ποιότητας δεδομένα, ενώ άλλες μπορεί να έχουν ελλείψεις ή ασυνέπειες. Για την Ελλάδα (ζώνη GR), τα δεδομένα είναι γενικά διαθέσιμα από το 2015 και μετά με καλή ποιότητα.

Οι κύριες προκλήσεις που αντιμετωπίζονται στην επεξεργασία των δεδομένων περιλαμβάνουν διαφορετικές χρονικές αναλύσεις (ωριαία, ημίωρη, τεταρτοωριαία) που απαιτούν ομογενοποίηση, ελλιπείς εγγραφές που απαιτούν συμπλήρωση ή αποκλεισμό, διαφορετικές ζώνες ώρας και μετάβαση θερινής/χειμερινής ώρας, και αλλαγές στη δομή των δεδομένων κατά τη διάρκεια του χρόνου.


% =============================================================================
\section{Αγωγός Δεδομένων (Data Pipeline)}
\label{sec:data-pipeline}

Για τη συστηματική συλλογή και ενημέρωση των δεδομένων από την πλατφόρμα ENTSO-E, αναπτύχθηκε ένας αυτοματοποιημένος αγωγός δεδομένων (data pipeline) σε Python. Ο αγωγός υλοποιεί ένα σύστημα σταδιακών ενημερώσεων (incremental updates) που επιτρέπει την αποτελεσματική ενημέρωση της βάσης δεδομένων χωρίς την ανάγκη πλήρους επαναφόρτωσης~\cite{kleppmann2017}.

\subsection{Αρχιτεκτονική Pipeline}
\label{subsec:pipeline-architecture}

Ο αγωγός δεδομένων αποτελείται από τα ακόλουθα βασικά στοιχεία:

\paragraph{Κεντρική Βιβλιοθήκη (helper\_functions.py)}
Περιέχει τις βασικές λειτουργίες που χρησιμοποιούνται από όλα τα scripts ενημέρωσης, συμπεριλαμβανομένων συναρτήσεων αυθεντικοποίησης με την πλατφόρμα ENTSO-E, ανακάλυψης και φιλτραρίσματος αρχείων, λήψης και επεξεργασίας CSV, μαζικής φόρτωσης στη βάση δεδομένων, και διαχείρισης αρχείων καταγραφής.

\paragraph{Scripts Ενημέρωσης (update\_*.py)}
Για κάθε τύπο δεδομένων υπάρχει ένα ξεχωριστό script ενημέρωσης που καθορίζει τις παραμέτρους ειδικά για αυτόν τον τύπο (όνομα πίνακα, φάκελος ENTSO-E, διαχωριστικό αρχείου, κωδικοποίηση κ.λπ.). Αυτή η αρχιτεκτονική επιτρέπει την εύκολη προσθήκη νέων τύπων δεδομένων και την ανεξάρτητη εκτέλεση κάθε ενημέρωσης.

\paragraph{Εργαλείο Δημιουργίας Scripts (generate\_incremental\_scripts.py)}
Αυτοματοποιεί τη δημιουργία των scripts ενημέρωσης με βάση πρότυπα, διασφαλίζοντας συνέπεια στη δομή και τη λειτουργικότητα.

\subsection{Σύστημα Σταδιακών Ενημερώσεων}
\label{subsec:incremental-updates}

Η κεντρική καινοτομία του αγωγού δεδομένων είναι το σύστημα σταδιακών ενημερώσεων που αποφεύγει την επαναφόρτωση ολόκληρου του συνόλου δεδομένων σε κάθε εκτέλεση. Η ροή εργασιών περιγράφεται ως εξής:

Αρχικά, το σύστημα εκτελεί αυθεντικοποίηση με την πλατφόρμα ENTSO-E χρησιμοποιώντας διαπιστευτήρια που αποθηκεύονται σε μεταβλητές περιβάλλοντος. Στη συνέχεια, συνδέεται μέσω SFTP και ανακαλύπτει τα διαθέσιμα αρχεία στον αντίστοιχο φάκελο, καταγράφοντας το όνομα και την ημερομηνία τελευταίας τροποποίησης κάθε αρχείου. Το σύστημα συγκρίνει τη λίστα με τον πίνακα \texttt{file\_update\_log} της βάσης δεδομένων και εντοπίζει τα νέα ή τροποποιημένα αρχεία. Τα αρχεία επεξεργάζονται σε παρτίδες (batches) καθορισμένου μεγέθους (προεπιλογή: 50 αρχεία) για αποτελεσματική χρήση μνήμης. Κάθε αρχείο CSV αναλύεται, καθαρίζεται (data cleaning) και μετασχηματίζεται στη μορφή του πίνακα προορισμού. Τέλος, τα δεδομένα φορτώνονται στη βάση χρησιμοποιώντας τη μέθοδο PostgreSQL COPY για μέγιστη απόδοση, και ο πίνακας \texttt{file\_update\_log} ενημερώνεται με τα στοιχεία των επεξεργασμένων αρχείων.

\subsection{Βελτιστοποίηση Απόδοσης}
\label{subsec:pipeline-performance}

Για τη βελτιστοποίηση της απόδοσης του αγωγού δεδομένων, υλοποιήθηκαν διάφορες τεχνικές. Η χρήση της μεθόδου PostgreSQL COPY αντί για πολλαπλές εντολές INSERT προσφέρει σημαντικά υψηλότερη απόδοση για μαζική φόρτωση δεδομένων. Το ρυθμιζόμενο μέγεθος παρτίδας (batch size) επιτρέπει την εξισορρόπηση μεταξύ χρήσης μνήμης και απόδοσης. Η επεξεργασία σε παρτίδες αντί της φόρτωσης ολόκληρου του αρχείου στη μνήμη επιτρέπει τη διαχείριση πολύ μεγάλων συνόλων δεδομένων. Ο μηχανισμός παρακολούθησης αρχείων μέσω του \texttt{file\_update\_log} αποφεύγει την περιττή επαναεπεξεργασία αρχείων.

\subsection{Χρονοπρογραμματισμός}
\label{subsec:scheduling}

Η αυτοματοποιημένη εκτέλεση του αγωγού δεδομένων επιτυγχάνεται μέσω GitHub Actions, με τις εργασίες να εκτελούνται σε καθημερινή βάση. Η ροή εργασιών περιλαμβάνει την εκτέλεση του \texttt{generate\_incremental\_scripts.py} για τη δημιουργία των τελευταίων scripts, ακολουθούμενη από την εκτέλεση κάθε script ενημέρωσης για τους διάφορους τύπους δεδομένων.


% =============================================================================
\section{Σχήμα Βάσης Δεδομένων}
\label{sec:database-schema}

Η βάση δεδομένων PostgreSQL οργανώνεται σε δύο κύρια σχήματα (schemas): το σχήμα \texttt{entsoe} που περιέχει τα ακατέργαστα δεδομένα από την πλατφόρμα ENTSO-E, και το σχήμα \texttt{simulations} που περιέχει τα αποτελέσματα των προσομοιώσεων.

\subsection{Σχήμα ENTSO-E}
\label{subsec:entsoe-schema}

Το σχήμα \texttt{entsoe} περιέχει τους πίνακες που αποθηκεύουν τα δεδομένα από την πλατφόρμα διαφάνειας. Οι κύριοι πίνακες παρουσιάζονται στον Πίνακα~\ref{tab:entsoe-tables}.

\begin{table}[htbp]
\centering
\caption{Κύριοι πίνακες του σχήματος \texttt{entsoe}}
\label{tab:entsoe-tables}
\begin{tabular}{@{}p{5.5cm}p{8cm}@{}}
\toprule
\textbf{Πίνακας} & \textbf{Περιγραφή} \\
\midrule
\texttt{production\_and\_generation\_units} & Στοιχεία μονάδων παραγωγής \\
\texttt{actual\_total\_load} & Ιστορικά δεδομένα κατανάλωσης \\
\texttt{actual\_generation\_output\_per\_generation\_unit} & Ιστορική παραγωγή ανά μονάδα \\
\texttt{generation\_forecasts\_for\_wind\_and\_solar} & Προβλέψεις παραγωγής ΑΠΕ \\
\texttt{offered\_transfer\_capacities\_implicit} & Διαθέσιμη ικανότητα μεταφοράς \\
\texttt{unavailability\_of\_production\_and\_generation\_units} & Μη διαθεσιμότητα μονάδων \\
\texttt{energy\_prices} & Τιμές εκκαθάρισης αγοράς \\
\texttt{physical\_flows} & Φυσικές ροές διασυνδέσεων \\
\texttt{file\_update\_log} & Καταγραφή επεξεργασμένων αρχείων \\
\bottomrule
\end{tabular}
\end{table}

\paragraph{Πίνακας production\_and\_generation\_units}
Ο πίνακας αυτός περιέχει τα στοιχεία των μονάδων παραγωγής με τα ακόλουθα βασικά πεδία: \texttt{generation\_unit\_code} (μοναδικός κωδικός μονάδας), \texttt{generation\_unit\_name} (όνομα μονάδας), \texttt{area\_map\_code} (κωδικός ζώνης προσφοράς, π.χ. ``GR''), \texttt{production\_unit\_type} (τύπος καυσίμου), \texttt{generation\_unit\_installed\_capacity\_mw} (εγκατεστημένη ισχύς), \texttt{production\_unit\_status} (κατάσταση λειτουργίας), \texttt{valid\_from} και \texttt{valid\_to} (περίοδος ισχύος των στοιχείων).

Η ερώτηση για την ανάκτηση διαθέσιμων μονάδων για μια συγκεκριμένη ζώνη και ημερομηνία φιλτράρει τις μονάδες με βάση την κατάσταση \texttt{COMMISSIONED}, τον τύπο περιοχής \texttt{BZN} ή \texttt{BZN/CTA}, και την περίοδο ισχύος.

\paragraph{Πίνακας actual\_total\_load}
Ο πίνακας αυτός περιέχει χρονοσειρές κατανάλωσης με τα πεδία: \texttt{datetime\_utc} (χρονοσφραγίδα σε UTC), \texttt{area\_map\_code} (ζώνη προσφοράς), \texttt{total\_load\_mw} (κατανάλωση σε MW), και \texttt{resolution} (χρονική ανάλυση).

\paragraph{Πίνακας actual\_generation\_output\_per\_generation\_unit}
Ο πίνακας αυτός περιέχει ιστορικά δεδομένα πραγματικής παραγωγής ανά μονάδα με τα πεδία: \texttt{generation\_unit\_code} (κωδικός μονάδας), \texttt{datetime\_utc} (χρονοσφραγίδα σε UTC), \texttt{actual\_generation\_output\_mw} (μετρούμενη παραγωγή σε MW), και \texttt{resolution} (χρονική ανάλυση, συνήθως 15 λεπτά ή 1 ώρα). Τα δεδομένα αυτά χρησιμοποιούνται για την εξαγωγή τεχνικών παραμέτρων (ρυθμοί μεταβολής, ελάχιστη ισχύς, χρόνοι λειτουργίας/αδράνειας) μέσω στατιστικής ανάλυσης, όπως περιγράφεται στην Ενότητα~\ref{subsec:parameter-inference}.

\paragraph{Πίνακας unavailability\_of\_production\_and\_generation\_units}
Ο πίνακας αυτός περιέχει πληροφορίες διακοπών με τα πεδία: \texttt{asset\_code} (αντιστοιχεί στο \texttt{generation\_unit\_code}), \texttt{start\_outage\_utc} και \texttt{end\_outage\_utc} (περίοδος διακοπής), \texttt{type} (``Planned'' ή ``Forced''), \texttt{status} (``Active'', ``Cancelled'', ``Withdrawn''), και \texttt{available\_capacity\_mw} (εναπομένουσα διαθέσιμη ισχύς).

\subsection{Σχήμα Προσομοιώσεων}
\label{subsec:simulations-schema}

Το σχήμα \texttt{simulations} περιέχει τους πίνακες για την αποθήκευση των αποτελεσμάτων των προσομοιώσεων. Οι κύριοι πίνακες παρουσιάζονται στον Πίνακα~\ref{tab:simulations-tables}.

\begin{table}[htbp]
\centering
\caption{Κύριοι πίνακες του σχήματος \texttt{simulations}}
\label{tab:simulations-tables}
\begin{tabular}{@{}p{5cm}p{8cm}@{}}
\toprule
\textbf{Πίνακας} & \textbf{Περιγραφή} \\
\midrule
\texttt{energy\_prices} & Υπολογισμένες τιμές εκκαθάρισης \\
\texttt{optimization\_runs} & Μεταδεδομένα εκτελέσεων βελτιστοποίησης \\
\texttt{transmission\_flows} & Διασυνοριακές ροές ισχύος \\
\texttt{generator\_inferred\_parameters} & Εξαχθείσες τεχνικές παράμετροι μονάδων \\
\texttt{uc\_results} & Αποτελέσματα Unit Commitment (σύνοψη) \\
\texttt{uc\_generation} & Παραγωγή/δέσμευση ανά μονάδα και περίοδο \\
\texttt{uc\_net\_demand} & Καθαρή ζήτηση ανά περίοδο \\
\bottomrule
\end{tabular}
\end{table}

\paragraph{Πίνακας energy\_prices}
Ο πίνακας αυτός αποθηκεύει τις υπολογισμένες τιμές εκκαθάρισης με τα πεδία: \texttt{bidding\_zone} (ζώνη προσφοράς), \texttt{date} (ημερομηνία), \texttt{hour} (ώρα 1-24), \texttt{price} (τιμή σε €/MWh), \texttt{created\_at} (χρονοσφραγίδα δημιουργίας), και \texttt{run\_id} (αναφορά στην εκτέλεση βελτιστοποίησης).

\paragraph{Πίνακας transmission\_flows}
Ο πίνακας αυτός αποθηκεύει τις υπολογισμένες διασυνοριακές ροές με τα πεδία: \texttt{source\_zone} και \texttt{target\_zone} (ζώνες προέλευσης και προορισμού), \texttt{date} (ημερομηνία), \texttt{hour} (ώρα), \texttt{flow\_mw} (ροή σε MW), και \texttt{run\_id}.

\paragraph{Πίνακας generator\_inferred\_parameters}
Ο πίνακας αυτός λειτουργεί ως κρυφή μνήμη (cache) για τις τεχνικές παραμέτρους μονάδων παραγωγής που εξάγονται από ιστορικά δεδομένα. Τα βασικά πεδία περιλαμβάνουν: \texttt{generator\_code} (κωδικός μονάδας), \texttt{ramp\_up} και \texttt{ramp\_down} (ρυθμοί μεταβολής ως κλάσμα της $P^{\max}$ ανά ώρα), \texttt{p\_min} (ελάχιστη ισχύς ως κλάσμα της $P^{\max}$), \texttt{min\_uptime} και \texttt{min\_downtime} (ελάχιστοι χρόνοι λειτουργίας/αδράνειας σε ώρες), και \texttt{inference\_date} (ημερομηνία εξαγωγής για έλεγχο λήξης). Η χρήση του πίνακα αυτού μειώνει τον χρόνο ανάκτησης παραμέτρων από περίπου 17 λεπτά σε 2 δευτερόλεπτα.

\paragraph{Πίνακες uc\_results, uc\_generation, uc\_net\_demand}
Οι τρεις αυτοί πίνακες αποτελούν τον μηχανισμό κρυφής μνήμης (caching) για τα αποτελέσματα του επιλυτή Unit Commitment. Ο πίνακας \texttt{uc\_results} αποθηκεύει τη σύνοψη κάθε εκτέλεσης με κλειδί \texttt{(bidding\_zone, market\_date, code\_version)} και περιλαμβάνει τα πεδία: \texttt{status} (κατάσταση τερματισμού), \texttt{solver} (επιλυτής), \texttt{resolution\_minutes} (χρονική ανάλυση), \texttt{total\_cost}, \texttt{production\_cost}, \texttt{startup\_cost}, \texttt{noload\_cost} (συνιστώσες κόστους), και αριθμούς εκκινήσεων (\texttt{hot/warm/cold\_startups}). Ο πίνακας \texttt{uc\_generation} αποθηκεύει τα αναλυτικά αποτελέσματα παραγωγής, δέσμευσης και εκκίνησης ανά μονάδα και περίοδο, ενώ ο πίνακας \texttt{uc\_net\_demand} αποθηκεύει την καθαρή ζήτηση και την παραγωγή ΑΠΕ ανά περίοδο. Η αποθήκευση ανέρχεται σε περίπου 195 KB ανά ζώνη/ημέρα.


% =============================================================================
\section{Οικοσύστημα Julia/JuMP}
\label{sec:julia-jump}

Για την υλοποίηση των αλγορίθμων βελτιστοποίησης επιλέχθηκε η γλώσσα προγραμματισμού Julia σε συνδυασμό με το πλαίσιο μοντελοποίησης JuMP~\cite{bezanson2017, jump2017}. Αυτή η επιλογή βασίζεται σε τεχνικά πλεονεκτήματα που καθιστούν το συγκεκριμένο οικοσύστημα ιδανικό για επιστημονικούς υπολογισμούς και προβλήματα βελτιστοποίησης.

\subsection{Γλώσσα Προγραμματισμού Julia}
\label{subsec:julia-language}

Η Julia είναι μια σύγχρονη γλώσσα προγραμματισμού υψηλού επιπέδου σχεδιασμένη για επιστημονικούς υπολογισμούς. Τα κύρια χαρακτηριστικά της που την καθιστούν κατάλληλη για την παρούσα εφαρμογή είναι τα ακόλουθα.

Η Julia επιτυγχάνει υψηλή απόδοση μέσω της μεταγλώττισης just-in-time (JIT) με χρήση του LLVM compiler framework. Ο κώδικας Julia μεταγλωττίζεται σε εγγενή κώδικα μηχανής, επιτυγχάνοντας ταχύτητες συγκρίσιμες με C και Fortran, χωρίς την ανάγκη για χειροκίνητη βελτιστοποίηση ή κλήση εξωτερικών βιβλιοθηκών.

Παρά την υψηλή απόδοση, η σύνταξη της Julia είναι εκφραστική και φιλική στον χρήστη, παρόμοια με Python ή MATLAB. Αυτό επιτρέπει την ταχεία ανάπτυξη πρωτοτύπων και την εύκολη ανάγνωση του κώδικα.

Η Julia διαθέτει ένα εξελιγμένο σύστημα τύπων με πολλαπλή αποστολή (multiple dispatch) που επιτρέπει τη δημιουργία γενικευμένου κώδικα που είναι ταυτόχρονα αποτελεσματικός και επαναχρησιμοποιήσιμος.

Η Julia παρέχει ενσωματωμένη υποστήριξη για παραλληλισμό και κατανεμημένους υπολογισμούς, διευκολύνοντας την κλιμάκωση σε πολλαπλούς πυρήνες ή μηχανές.

\subsection{Πλαίσιο Μοντελοποίησης JuMP}
\label{subsec:jump-framework}

Το JuMP (Julia for Mathematical Programming) είναι ένα domain-specific language (DSL) ενσωματωμένο στη Julia για τη διατύπωση και επίλυση προβλημάτων μαθηματικού προγραμματισμού. Αποτελεί ένα από τα πιο δημοφιλή και ώριμα πλαίσια μοντελοποίησης βελτιστοποίησης στον ακαδημαϊκό και βιομηχανικό χώρο~\cite{jump2017}.

Τα πλεονεκτήματα του JuMP περιλαμβάνουν μια αλγεβρική μοντελοποίηση που επιτρέπει τη διατύπωση προβλημάτων βελτιστοποίησης με τρόπο που μοιάζει με τη μαθηματική τους διατύπωση. Η σύνταξη για μεταβλητές, περιορισμούς και αντικειμενικές συναρτήσεις είναι διαισθητική και εκφραστική. Το JuMP υποστηρίζει πολλαπλούς επιλυτές (solvers) μέσω μιας ενιαίας διεπαφής, επιτρέποντας την εύκολη εναλλαγή μεταξύ επιλυτών χωρίς αλλαγές στον κώδικα μοντελοποίησης. Η απόδοση του JuMP είναι συγκρίσιμη ή καλύτερη από εμπορικά εργαλεία μοντελοποίησης όπως το AMPL και το GAMS. Τέλος, το JuMP είναι εντελώς ελεύθερο και ανοιχτού κώδικα, διατίθεται υπό την άδεια MIT.

Ένα τυπικό παράδειγμα χρήσης του JuMP για τη δημιουργία ενός μοντέλου βελτιστοποίησης φαίνεται στο ακόλουθο απόσπασμα κώδικα:

\begin{samepage}
\begin{minted}[fontsize=\small,linenos,frame=lines]{julia}
using JuMP, HiGHS

# Δημιουργία μοντέλου με τον επιλυτή HiGHS
model = Model(HiGHS.Optimizer)

# Ορισμός μεταβλητών
@variable(model, x >= 0)
@variable(model, y >= 0)

# Ορισμός περιορισμών
@constraint(model, x + 2y <= 10)
@constraint(model, 3x + y <= 15)

# Ορισμός αντικειμενικής συνάρτησης
@objective(model, Max, 2x + 3y)

# Επίλυση
optimize!(model)
\end{minted}
\end{samepage}

\subsection{Επιλυτής HiGHS}
\label{subsec:highs-solver}

Ως κύριος επιλυτής για τα προβλήματα βελτιστοποίησης χρησιμοποιείται ο HiGHS, ένας σύγχρονος επιλυτής ανοιχτού κώδικα για προβλήματα γραμμικού προγραμματισμού (LP), μικτού ακέραιου προγραμματισμού (MIP) και τετραγωνικού προγραμματισμού (QP)~\cite{highs2024}.

Ο HiGHS αναπτύσσεται στο Πανεπιστήμιο του Εδιμβούργου και αποτελεί έναν από τους ταχύτερους ανοιχτού κώδικα επιλυτές. Περιλαμβάνει υλοποιήσεις του αλγορίθμου dual revised simplex με εκτεταμένη παραλληλοποίηση, μεθόδων εσωτερικού σημείου (interior point methods), και αλγορίθμων branch-and-bound για MIP με cutting planes και heuristics.

Η διεπαφή με τη Julia παρέχεται μέσω του πακέτου HiGHS.jl, το οποίο ενσωματώνεται απρόσκοπτα με το JuMP. Η επιλογή του HiGHS ως προεπιλεγμένου επιλυτή βασίζεται στην άδεια ανοιχτού κώδικα (MIT) που επιτρέπει ελεύθερη χρήση, στην υψηλή απόδοση συγκρίσιμη με εμπορικούς επιλυτές, στην ενεργή ανάπτυξη και υποστήριξη, και στη συμβατότητα με όλα τα απαιτούμενα είδη προβλημάτων (LP, MILP).

Το σύστημα υποστηρίζει επίσης εμπορικούς επιλυτές όπως ο Gurobi και ο CPLEX για περιπτώσεις που απαιτείται ακόμη υψηλότερη απόδοση, με την εναλλαγή να γίνεται μέσω μιας απλής παραμέτρου.

\subsection{Βασικά Πακέτα Julia}
\label{subsec:julia-packages}

Εκτός από το JuMP και τον HiGHS, το σύστημα χρησιμοποιεί μια σειρά από πακέτα Julia για διάφορες λειτουργίες. Το πακέτο DataFrames.jl χρησιμοποιείται για τη διαχείριση και επεξεργασία δεδομένων σε μορφή πινάκων, με λειτουργικότητα παρόμοια με τα pandas της Python. Το πακέτο LibPQ.jl παρέχει σύνδεση με τη βάση δεδομένων PostgreSQL μέσω της βιβλιοθήκης libpq, υποστηρίζοντας παραμετροποιημένα ερωτήματα και διαχείριση συνδέσεων. Το πακέτο CSV.jl χρησιμοποιείται για την ανάγνωση και εγγραφή αρχείων CSV με υψηλή απόδοση. Το πακέτο Dates.jl, που αποτελεί μέρος της standard library, χρησιμοποιείται για τη διαχείριση ημερομηνιών και χρόνων. Το πακέτο Statistics.jl, επίσης μέρος της standard library, χρησιμοποιείται για στατιστικούς υπολογισμούς---κυρίως τη συνάρτηση \texttt{quantile} για τον υπολογισμό εκατοστημορίων κατά την εξαγωγή τεχνικών παραμέτρων μονάδων από ιστορικά δεδομένα. Τέλος, το πακέτο DotEnv.jl χρησιμοποιείται για τη φόρτωση μεταβλητών περιβάλλοντος από αρχεία \texttt{.env} για τη διαχείριση διαπιστευτηρίων.


% =============================================================================
\section{Αρχιτεκτονική Εφαρμογής}
\label{sec:application-architecture}

Η εφαρμογή Julia οργανώνεται ως ένα ενιαίο module ονόματι \texttt{Euphemia} που περιέχει όλη τη λειτουργικότητα του συστήματος. Η δομή του module ακολουθεί τις βέλτιστες πρακτικές της Julia για οργάνωση κώδικα και επαναχρησιμοποίηση~\cite{julia-docs}.

\subsection{Δομή Module}
\label{subsec:module-structure}

Το module \texttt{Euphemia} αποτελείται από τα ακόλουθα υπο-modules και αρχεία, όπως φαίνεται στο Σχήμα~\ref{fig:module-structure}:

\begin{figure}[htbp]
\centering
\begin{verbatim}
Euphemia.jl (Main Module)
├── dbutils.jl                 # Database utilities
├── Generators.jl              # Generator data model
├── Loads.jl                   # Load data model
├── Renewables.jl              # RES forecasts
├── FuelTypeParameters.jl      # Fuel-specific parameters
├── MarketOrders.jl            # Order types
├── UnitCommitment.jl          # UC solver
├── BiddingStrategy.jl         # UC → Market orders
├── Network.jl                 # ATC constraints
├── MPCC.jl                    # Market clearing
└── AlternativeOrderBook.jl    # Fast order generation
\end{verbatim}
\caption{Δομή του module \texttt{Euphemia}.}
\label{fig:module-structure}
\end{figure}

\paragraph{Κεντρικό Module (Euphemia.jl)}
Αποτελεί το σημείο εισόδου της εφαρμογής. Διαχειρίζεται τις εξαρτήσεις, συντονίζει τα υπο-modules, και εξάγει τη δημόσια διεπαφή (public API). Περιλαμβάνει επίσης τη λογική επιλογής επιλυτή και την κεντρική ροή εργασιών.

\paragraph{Βοηθητικά Βάσης Δεδομένων (dbutils.jl)}
Περιέχει συναρτήσεις για τη σύνδεση με τη βάση δεδομένων PostgreSQL και την εκτέλεση ερωτημάτων. Υλοποιεί το pattern \texttt{withdb} για ασφαλή διαχείριση συνδέσεων με αυτόματο κλείσιμο.

\paragraph{Μοντέλα Δεδομένων (Generators.jl, Loads.jl, Renewables.jl)}
Ορίζουν τις δομές δεδομένων για τις οντότητες του συστήματος και παρέχουν συναρτήσεις για την ανάκτησή τους από τη βάση δεδομένων. Η δομή \texttt{Generator} περιλαμβάνει πεδία για τον κωδικό, το όνομα, τη ζώνη προσφοράς, τον τύπο καυσίμου, τα όρια ισχύος, το οριακό κόστος, και προαιρετικούς τεχνικούς περιορισμούς (ρυθμοί μεταβολής, ελάχιστοι χρόνοι λειτουργίας/αδράνειας) που εξάγονται αυτόματα από ιστορικά δεδομένα. Η δομή \texttt{InitialConditions} αναπαριστά την αρχική κατάσταση κάθε μονάδας στη χρονική στιγμή $t=0$ (δέσμευση, ισχύς εξόδου, θερμοκρασιακή κατάσταση) και εξάγεται από ιστορικά δεδομένα 72 ωρών μέσω μαζικού ερωτήματος. Το αρχείο \texttt{Generators.jl} περιλαμβάνει τις συναρτήσεις εξαγωγής παραμέτρων (\texttt{infer\_ramp\_rates}, \texttt{infer\_p\_min}, \texttt{infer\_uptime\_downtime}), τη συνάρτηση \texttt{get\_initial\_conditions} για αυτόματο προσδιορισμό αρχικών συνθηκών, και τον μηχανισμό αποθήκευσης σε κρυφή μνήμη βάσης δεδομένων.

\paragraph{Τύποι Εντολών (MarketOrders.jl)}
Ορίζει τις δομές δεδομένων για τους διάφορους τύπους εντολών αγοράς, συμπεριλαμβανομένων των \texttt{SimpleOrder}, \texttt{BlockOrder}, \texttt{LinkedBlockOrder} κ.λπ.

\paragraph{Unit Commitment (UnitCommitment.jl)}
Υλοποιεί τον επιλυτή για το πρόβλημα Unit Commitment χρησιμοποιώντας JuMP/HiGHS. Περιλαμβάνει τη διατύπωση του μοντέλου με τριμερή αντικειμενική συνάρτηση (κόστος παραγωγής, κόστος εκκίνησης ανά θερμοκρασιακή κατάσταση, κόστος αδράνειας), την ενσωμάτωση αρχικών συνθηκών, ρυθμιζόμενες παραμέτρους επιλυτή (MIP gap, χρονικό όριο μέσω MOI), καθώς και μηχανισμό αποθήκευσης αποτελεσμάτων σε κρυφή μνήμη βάσης δεδομένων.

\paragraph{Στρατηγική Προσφορών (BiddingStrategy.jl)}
Υλοποιεί τη μετατροπή των αποτελεσμάτων του Unit Commitment σε εντολές αγοράς σύμφωνα με τις τρεις στρατηγικές που περιγράφηκαν στο Κεφάλαιο~\ref{sec:bidding-strategies}.

\paragraph{Δίκτυο (Network.jl)}
Διαχειρίζεται τους περιορισμούς δικτύου και τη διαθέσιμη ικανότητα μεταφοράς (ATC) μεταξύ ζωνών προσφοράς. Παρέχει λειτουργίες για την ανάκτηση δεδομένων ATC από το ENTSO-E και την προσθήκη περιορισμών στο μοντέλο βελτιστοποίησης.

\paragraph{Εκκαθάριση Αγοράς (MPCC.jl)}
Υλοποιεί τον αλγόριθμο εκκαθάρισης της αγοράς με τη μέθοδο MPCC. Περιλαμβάνει τη δομή \texttt{MPCCOrderBook} για τη διαχείριση των εντολών και τη συνάρτηση \texttt{solve\_mpcc\_market\_clearing} για την επίλυση.

\subsection{Ροή Εργασιών}
\label{subsec:workflow}

Η κεντρική ροή εργασιών του συστήματος για τον υπολογισμό τιμών αγοράς περιλαμβάνει τα ακόλουθα βήματα. Αρχικά, γίνεται φόρτωση δεδομένων από τη βάση, συμπεριλαμβανομένων των διαθέσιμων μονάδων παραγωγής (με φιλτράρισμα μη διαθεσιμοτήτων), της ζήτησης για την επιλεγμένη ημερομηνία, και των προβλέψεων παραγωγής ΑΠΕ. Στη συνέχεια, επιλύεται το πρόβλημα Unit Commitment για τον προσδιορισμό του βέλτιστου προγράμματος λειτουργίας. Ακολουθεί η μετατροπή της λύσης του UC σε εντολές αγοράς σύμφωνα με την επιλεγμένη στρατηγική. Εκτελείται η εκκαθάριση της αγοράς μέσω του αλγορίθμου MPCC, που καθορίζει τις τιμές ισορροπίας. Τέλος, τα αποτελέσματα αποθηκεύονται στη βάση δεδομένων.

Η κύρια συνάρτηση για την εκτέλεση αυτής της ροής είναι η \texttt{generate\_energy\_prices}:

\begin{samepage}
\begin{minted}[fontsize=\small,linenos,frame=lines]{julia}
# Παράδειγμα χρήσης για μονοζωνική εκκαθάριση
prices = generate_energy_prices("GR", Date(2024, 6, 15);
    order_method = :uc_based,
    markup_factor = 1.1,
    save_to_db = true
)

# Πολυζωνική εκκαθάριση με περιορισμούς ATC
result = run_multi_zone_market_clearing(Date(2024, 6, 15);
    zones = ["GR", "BG", "RO"],
    order_method = :alternative,
    save_to_db = true
)
\end{minted}
\end{samepage}

\subsection{Διαχείριση Επιλυτών}
\label{subsec:solver-management}

Το σύστημα υποστηρίζει πολλαπλούς επιλυτές μέσω ενός μηχανισμού αυτόματης επιλογής και caching. Η συνάρτηση \texttt{select\_solver} ελέγχει τη διαθεσιμότητα των εγκατεστημένων επιλυτών (HiGHS, Gurobi, CPLEX) και επιλέγει τον κατάλληλο βάσει προτίμησης ή διαθεσιμότητας.

Για τον επιλυτή Gurobi, ο οποίος απαιτεί άδεια χρήσης με σημαντικό κόστος αρχικοποίησης, υλοποιήθηκε ένας μηχανισμός caching του περιβάλλοντος (environment) που αποφεύγει την επανάληψη της διαδικασίας αυθεντικοποίησης σε διαδοχικές κλήσεις.

\begin{samepage}
\begin{minted}[fontsize=\small,linenos,frame=lines]{julia}
# Αυτόματη επιλογή επιλυτή
optimizer, solver_name = select_solver("auto")

# Προτίμηση συγκεκριμένου επιλυτή
optimizer, solver_name = select_solver("gurobi")

# Δημιουργία μοντέλου με τον επιλεγμένο επιλυτή
model = Model(optimizer)
\end{minted}
\end{samepage}

% =============================================================================

\chapter{Αρχιτεκτονική Συστήματος και Δεδομένα}
\label{ch:system-architecture}

Στο παρόν κεφάλαιο περιγράφεται η αρχιτεκτονική του συστήματος που αναπτύχθηκε για την προσομοίωση της Προημερήσιας Αγοράς ηλεκτρικής ενέργειας. Η αρχιτεκτονική βασίζεται σε τρεις πυλώνες: την υποδομή δεδομένων που συλλέγει και αποθηκεύει δεδομένα από την πλατφόρμα διαφάνειας του ENTSO-E, τη βάση δεδομένων PostgreSQL που οργανώνει τα δεδομένα σε κατάλληλο σχήμα, και τον πυρήνα υπολογισμών σε Julia/JuMP που υλοποιεί τους αλγορίθμους βελτιστοποίησης.

Η γενική ροή δεδομένων του συστήματος απεικονίζεται στο Σχήμα~\ref{fig:system-architecture}. Τα δεδομένα συλλέγονται από την πλατφόρμα ENTSO-E μέσω ενός αυτοματοποιημένου pipeline σε Python, αποθηκεύονται στη βάση δεδομένων PostgreSQL, επεξεργάζονται από τους αλγορίθμους Julia/JuMP, και τα αποτελέσματα αποθηκεύονται πίσω στη βάση για ανάλυση και οπτικοποίηση.

% TODO: Εισαγωγή σχήματος αρχιτεκτονικής
\begin{figure}[htbp]
\centering
% \includegraphics[width=0.95\textwidth]{figures/system_architecture.pdf}
\fbox{\parbox{0.9\textwidth}{\centering\vspace{3cm}
\texttt{ENTSO-E Transparency Platform} \\
$\downarrow$ \textit{(Python pipeline)} \\
\texttt{PostgreSQL Database} \\
$\downarrow$ \textit{(Julia/JuMP)} \\
\texttt{Euphemia Engine} \\
$\downarrow$ \\
\texttt{simulations.energy\_prices}
\vspace{1cm}}}
\caption{Γενική αρχιτεκτονική του συστήματος προσομοίωσης αγοράς ηλεκτρικής ενέργειας.}
\label{fig:system-architecture}
\end{figure}


% =============================================================================
\section{Πηγές Δεδομένων}
\label{sec:data-sources}

Η κύρια πηγή δεδομένων για το σύστημα είναι η Πλατφόρμα Διαφάνειας του ENTSO-E (ENTSO-E Transparency Platform), η οποία παρέχει δημοσίως διαθέσιμα δεδομένα για τις ευρωπαϊκές αγορές ηλεκτρικής ενέργειας σύμφωνα με τον Κανονισμό (ΕΕ) 543/2013~\cite{entsoe-transparency, eu-regulation-543}.

\subsection{Πλατφόρμα Διαφάνειας ENTSO-E}
\label{subsec:entsoe-platform}

Το ENTSO-E (European Network of Transmission System Operators for Electricity) είναι ο οργανισμός που συντονίζει τους Διαχειριστές Συστημάτων Μεταφοράς (TSOs) της Ευρώπης. Η Πλατφόρμα Διαφάνειας που διατηρεί παρέχει πρόσβαση σε δεδομένα που καλύπτουν τη λειτουργία των ευρωπαϊκών αγορών ηλεκτρικής ενέργειας, συμπεριλαμβανομένων δεδομένων παραγωγής, κατανάλωσης, τιμών, διασυνοριακών ροών και διαθεσιμότητας μονάδων.

Η πλατφόρμα προσφέρει δύο κύριους τρόπους πρόσβασης. Ο πρώτος είναι το RESTful API, το οποίο επιτρέπει προγραμματιστική πρόσβαση στα δεδομένα μέσω HTTP αιτημάτων. Το API χρησιμοποιεί μορφή XML για τις απαντήσεις και απαιτεί κλειδί αυθεντικοποίησης (API key) που παρέχεται δωρεάν μετά από εγγραφή. Ο δεύτερος τρόπος είναι η SFTP πρόσβαση, η οποία παρέχει πρόσβαση σε προ-επεξεργασμένα αρχεία CSV ανά τύπο δεδομένων. Αυτή η μέθοδος είναι αποτελεσματικότερη για μαζική λήψη δεδομένων και χρησιμοποιείται στην παρούσα υλοποίηση.

\subsection{Τύποι Δεδομένων}
\label{subsec:data-types}

Για την προσομοίωση της Προημερήσιας Αγοράς, το σύστημα συλλέγει και επεξεργάζεται τους ακόλουθους τύπους δεδομένων:

\paragraph{Μονάδες Παραγωγής (Production and Generation Units)}
Ο πίνακας \texttt{production\_and\_generation\_units} περιέχει πληροφορίες για τις μονάδες παραγωγής ηλεκτρικής ενέργειας σε κάθε ζώνη προσφοράς. Τα δεδομένα περιλαμβάνουν τον κωδικό και το όνομα της μονάδας, τη ζώνη προσφοράς (bidding zone), τον τύπο καυσίμου (π.χ. φυσικό αέριο, λιγνίτης, υδροηλεκτρικά, αιολικά), την εγκατεστημένη ισχύ σε MW, και την κατάσταση λειτουργίας (commissioned/decommissioned). Αυτά τα δεδομένα χρησιμοποιούνται για τη δημιουργία του συνόλου των διαθέσιμων μονάδων στο πρόβλημα Unit Commitment.

\paragraph{Πραγματικό Φορτίο (Actual Total Load)}
Ο πίνακας \texttt{actual\_total\_load} περιέχει ιστορικά δεδομένα κατανάλωσης ηλεκτρικής ενέργειας ανά ζώνη προσφοράς με ωριαία ή τεταρτοωριαία ανάλυση. Η ζήτηση αποτελεί βασική παράμετρο εισόδου τόσο στο πρόβλημα Unit Commitment όσο και στην εκκαθάριση της αγοράς.

\paragraph{Προβλέψεις Παραγωγής ΑΠΕ (Generation Forecasts for Wind and Solar)}
Ο πίνακας \texttt{generation\_forecasts\_for\_wind\_and\_solar} περιέχει προβλέψεις παραγωγής από αιολικά και φωτοβολταϊκά συστήματα. Δεδομένου ότι η παραγωγή από ΑΠΕ θεωρείται μη κατανεμόμενη (non-dispatchable), οι προβλέψεις αφαιρούνται από τη συνολική ζήτηση για τον υπολογισμό της καθαρής ζήτησης (net demand) που πρέπει να καλυφθεί από τις θερμικές μονάδες.

\paragraph{Διαθέσιμη Ικανότητα Μεταφοράς (Offered Transfer Capacities)}
Ο πίνακας \texttt{offered\_transfer\_capacities\_implicit} περιέχει δεδομένα για τη διαθέσιμη ικανότητα μεταφοράς (ATC) μεταξύ γειτονικών ζωνών προσφοράς. Τα δεδομένα περιλαμβάνουν τη ζώνη προέλευσης και προορισμού, τη χρονική περίοδο, και τη διαθέσιμη ικανότητα σε MW. Αυτά τα δεδομένα είναι απαραίτητα για την πολυζωνική εκκαθάριση της αγοράς με περιορισμούς δικτύου.

\paragraph{Μη Διαθεσιμότητα Μονάδων (Unavailability of Production Units)}
Ο πίνακας \texttt{unavailability\_of\_production\_and\_generation\_units} περιέχει πληροφορίες για προγραμματισμένες ή απρόβλεπτες διακοπές λειτουργίας μονάδων παραγωγής. Κάθε εγγραφή περιλαμβάνει τον κωδικό της μονάδας, την περίοδο μη διαθεσιμότητας, τον τύπο (προγραμματισμένη ή έκτακτη), την κατάσταση (ενεργή, ακυρωμένη), και την εναπομένουσα διαθέσιμη ισχύ. Αυτά τα δεδομένα χρησιμοποιούνται για τον αποκλεισμό ή τη μείωση της διαθέσιμης ισχύος μονάδων κατά την επίλυση του Unit Commitment.

\paragraph{Πραγματική Παραγωγή ανά Μονάδα (Actual Generation Output)}
Ο πίνακας \texttt{actual\_generation\_output\_per\_generation\_unit} περιέχει ιστορικά δεδομένα πραγματικής παραγωγής ανά μονάδα, με χρονική ανάλυση 15 λεπτών ή 1 ώρας. Κάθε εγγραφή περιλαμβάνει τον κωδικό της μονάδας, τη χρονοσφραγίδα (UTC), τον κωδικό ανάλυσης (\texttt{PT15M}, \texttt{PT60M}) και τη μετρούμενη παραγωγή σε MW. Τα δεδομένα αυτά αξιοποιούνται για τη στατιστική εξαγωγή τεχνικών παραμέτρων (ρυθμοί μεταβολής, ελάχιστη ισχύς, χρόνοι λειτουργίας/αδράνειας) μέσω ανάλυσης ιστορικού παραθύρου 3 μηνών.

\paragraph{Τιμές Ενέργειας (Energy Prices)}
Ο πίνακας \texttt{energy\_prices} περιέχει τις πραγματικές τιμές εκκαθάρισης της Προημερήσιας Αγοράς ανά ζώνη προσφοράς. Αυτά τα δεδομένα χρησιμοποιούνται για την επικύρωση (validation) των αποτελεσμάτων του συστήματος προσομοίωσης.

\subsection{Διαθεσιμότητα και Ποιότητα Δεδομένων}
\label{subsec:data-quality}

Η ποιότητα και η διαθεσιμότητα των δεδομένων του ENTSO-E ποικίλλει ανάλογα με τη ζώνη προσφοράς και τον τύπο δεδομένων. Ορισμένες ζώνες (π.χ. Γερμανία, Γαλλία, Σκανδιναβικές χώρες) διαθέτουν πλήρη και υψηλής ποιότητας δεδομένα, ενώ άλλες μπορεί να έχουν ελλείψεις ή ασυνέπειες. Για την Ελλάδα (ζώνη GR), τα δεδομένα είναι γενικά διαθέσιμα από το 2015 και μετά με καλή ποιότητα.

Οι κύριες προκλήσεις που αντιμετωπίζονται στην επεξεργασία των δεδομένων περιλαμβάνουν διαφορετικές χρονικές αναλύσεις (ωριαία, ημίωρη, τεταρτοωριαία) που απαιτούν ομογενοποίηση, ελλιπείς εγγραφές που απαιτούν συμπλήρωση ή αποκλεισμό, διαφορετικές ζώνες ώρας και μετάβαση θερινής/χειμερινής ώρας, και αλλαγές στη δομή των δεδομένων κατά τη διάρκεια του χρόνου.


% =============================================================================
\section{Αγωγός Δεδομένων (Data Pipeline)}
\label{sec:data-pipeline}

Για τη συστηματική συλλογή και ενημέρωση των δεδομένων από την πλατφόρμα ENTSO-E, αναπτύχθηκε ένας αυτοματοποιημένος αγωγός δεδομένων (data pipeline) σε Python. Ο αγωγός υλοποιεί ένα σύστημα σταδιακών ενημερώσεων (incremental updates) που επιτρέπει την αποτελεσματική ενημέρωση της βάσης δεδομένων χωρίς την ανάγκη πλήρους επαναφόρτωσης~\cite{kleppmann2017}.

\subsection{Αρχιτεκτονική Pipeline}
\label{subsec:pipeline-architecture}

Ο αγωγός δεδομένων αποτελείται από τα ακόλουθα βασικά στοιχεία:

\paragraph{Κεντρική Βιβλιοθήκη (helper\_functions.py)}
Περιέχει τις βασικές λειτουργίες που χρησιμοποιούνται από όλα τα scripts ενημέρωσης, συμπεριλαμβανομένων συναρτήσεων αυθεντικοποίησης με την πλατφόρμα ENTSO-E, ανακάλυψης και φιλτραρίσματος αρχείων, λήψης και επεξεργασίας CSV, μαζικής φόρτωσης στη βάση δεδομένων, και διαχείρισης αρχείων καταγραφής.

\paragraph{Scripts Ενημέρωσης (update\_*.py)}
Για κάθε τύπο δεδομένων υπάρχει ένα ξεχωριστό script ενημέρωσης που καθορίζει τις παραμέτρους ειδικά για αυτόν τον τύπο (όνομα πίνακα, φάκελος ENTSO-E, διαχωριστικό αρχείου, κωδικοποίηση κ.λπ.). Αυτή η αρχιτεκτονική επιτρέπει την εύκολη προσθήκη νέων τύπων δεδομένων και την ανεξάρτητη εκτέλεση κάθε ενημέρωσης.

\paragraph{Εργαλείο Δημιουργίας Scripts (generate\_incremental\_scripts.py)}
Αυτοματοποιεί τη δημιουργία των scripts ενημέρωσης με βάση πρότυπα, διασφαλίζοντας συνέπεια στη δομή και τη λειτουργικότητα.

\subsection{Σύστημα Σταδιακών Ενημερώσεων}
\label{subsec:incremental-updates}

Η κεντρική καινοτομία του αγωγού δεδομένων είναι το σύστημα σταδιακών ενημερώσεων που αποφεύγει την επαναφόρτωση ολόκληρου του συνόλου δεδομένων σε κάθε εκτέλεση. Η ροή εργασιών περιγράφεται ως εξής:

Αρχικά, το σύστημα εκτελεί αυθεντικοποίηση με την πλατφόρμα ENTSO-E χρησιμοποιώντας διαπιστευτήρια που αποθηκεύονται σε μεταβλητές περιβάλλοντος. Στη συνέχεια, συνδέεται μέσω SFTP και ανακαλύπτει τα διαθέσιμα αρχεία στον αντίστοιχο φάκελο, καταγράφοντας το όνομα και την ημερομηνία τελευταίας τροποποίησης κάθε αρχείου. Το σύστημα συγκρίνει τη λίστα με τον πίνακα \texttt{file\_update\_log} της βάσης δεδομένων και εντοπίζει τα νέα ή τροποποιημένα αρχεία. Τα αρχεία επεξεργάζονται σε παρτίδες (batches) καθορισμένου μεγέθους (προεπιλογή: 50 αρχεία) για αποτελεσματική χρήση μνήμης. Κάθε αρχείο CSV αναλύεται, καθαρίζεται (data cleaning) και μετασχηματίζεται στη μορφή του πίνακα προορισμού. Τέλος, τα δεδομένα φορτώνονται στη βάση χρησιμοποιώντας τη μέθοδο PostgreSQL COPY για μέγιστη απόδοση, και ο πίνακας \texttt{file\_update\_log} ενημερώνεται με τα στοιχεία των επεξεργασμένων αρχείων.

\subsection{Βελτιστοποίηση Απόδοσης}
\label{subsec:pipeline-performance}

Για τη βελτιστοποίηση της απόδοσης του αγωγού δεδομένων, υλοποιήθηκαν διάφορες τεχνικές. Η χρήση της μεθόδου PostgreSQL COPY αντί για πολλαπλές εντολές INSERT προσφέρει σημαντικά υψηλότερη απόδοση για μαζική φόρτωση δεδομένων. Το ρυθμιζόμενο μέγεθος παρτίδας (batch size) επιτρέπει την εξισορρόπηση μεταξύ χρήσης μνήμης και απόδοσης. Η επεξεργασία σε παρτίδες αντί της φόρτωσης ολόκληρου του αρχείου στη μνήμη επιτρέπει τη διαχείριση πολύ μεγάλων συνόλων δεδομένων. Ο μηχανισμός παρακολούθησης αρχείων μέσω του \texttt{file\_update\_log} αποφεύγει την περιττή επαναεπεξεργασία αρχείων.

\subsection{Χρονοπρογραμματισμός}
\label{subsec:scheduling}

Η αυτοματοποιημένη εκτέλεση του αγωγού δεδομένων επιτυγχάνεται μέσω GitHub Actions, με τις εργασίες να εκτελούνται σε καθημερινή βάση. Η ροή εργασιών περιλαμβάνει την εκτέλεση του \texttt{generate\_incremental\_scripts.py} για τη δημιουργία των τελευταίων scripts, ακολουθούμενη από την εκτέλεση κάθε script ενημέρωσης για τους διάφορους τύπους δεδομένων.


% =============================================================================
\section{Σχήμα Βάσης Δεδομένων}
\label{sec:database-schema}

Η βάση δεδομένων PostgreSQL οργανώνεται σε δύο κύρια σχήματα (schemas): το σχήμα \texttt{entsoe} που περιέχει τα ακατέργαστα δεδομένα από την πλατφόρμα ENTSO-E, και το σχήμα \texttt{simulations} που περιέχει τα αποτελέσματα των προσομοιώσεων.

\subsection{Σχήμα ENTSO-E}
\label{subsec:entsoe-schema}

Το σχήμα \texttt{entsoe} περιέχει τους πίνακες που αποθηκεύουν τα δεδομένα από την πλατφόρμα διαφάνειας. Οι κύριοι πίνακες παρουσιάζονται στον Πίνακα~\ref{tab:entsoe-tables}.

\begin{table}[htbp]
\centering
\caption{Κύριοι πίνακες του σχήματος \texttt{entsoe}}
\label{tab:entsoe-tables}
\begin{tabular}{@{}p{5.5cm}p{8cm}@{}}
\toprule
\textbf{Πίνακας} & \textbf{Περιγραφή} \\
\midrule
\texttt{production\_and\_generation\_units} & Στοιχεία μονάδων παραγωγής \\
\texttt{actual\_total\_load} & Ιστορικά δεδομένα κατανάλωσης \\
\texttt{actual\_generation\_output\_per\_generation\_unit} & Ιστορική παραγωγή ανά μονάδα \\
\texttt{generation\_forecasts\_for\_wind\_and\_solar} & Προβλέψεις παραγωγής ΑΠΕ \\
\texttt{offered\_transfer\_capacities\_implicit} & Διαθέσιμη ικανότητα μεταφοράς \\
\texttt{unavailability\_of\_production\_and\_generation\_units} & Μη διαθεσιμότητα μονάδων \\
\texttt{energy\_prices} & Τιμές εκκαθάρισης αγοράς \\
\texttt{physical\_flows} & Φυσικές ροές διασυνδέσεων \\
\texttt{file\_update\_log} & Καταγραφή επεξεργασμένων αρχείων \\
\bottomrule
\end{tabular}
\end{table}

\paragraph{Πίνακας production\_and\_generation\_units}
Ο πίνακας αυτός περιέχει τα στοιχεία των μονάδων παραγωγής με τα ακόλουθα βασικά πεδία: \texttt{generation\_unit\_code} (μοναδικός κωδικός μονάδας), \texttt{generation\_unit\_name} (όνομα μονάδας), \texttt{area\_map\_code} (κωδικός ζώνης προσφοράς, π.χ. ``GR''), \texttt{production\_unit\_type} (τύπος καυσίμου), \texttt{generation\_unit\_installed\_capacity\_mw} (εγκατεστημένη ισχύς), \texttt{production\_unit\_status} (κατάσταση λειτουργίας), \texttt{valid\_from} και \texttt{valid\_to} (περίοδος ισχύος των στοιχείων).

Η ερώτηση για την ανάκτηση διαθέσιμων μονάδων για μια συγκεκριμένη ζώνη και ημερομηνία φιλτράρει τις μονάδες με βάση την κατάσταση \texttt{COMMISSIONED}, τον τύπο περιοχής \texttt{BZN} ή \texttt{BZN/CTA}, και την περίοδο ισχύος.

\paragraph{Πίνακας actual\_total\_load}
Ο πίνακας αυτός περιέχει χρονοσειρές κατανάλωσης με τα πεδία: \texttt{datetime\_utc} (χρονοσφραγίδα σε UTC), \texttt{area\_map\_code} (ζώνη προσφοράς), \texttt{total\_load\_mw} (κατανάλωση σε MW), και \texttt{resolution} (χρονική ανάλυση).

\paragraph{Πίνακας actual\_generation\_output\_per\_generation\_unit}
Ο πίνακας αυτός περιέχει ιστορικά δεδομένα πραγματικής παραγωγής ανά μονάδα με τα πεδία: \texttt{generation\_unit\_code} (κωδικός μονάδας), \texttt{datetime\_utc} (χρονοσφραγίδα σε UTC), \texttt{actual\_generation\_output\_mw} (μετρούμενη παραγωγή σε MW), και \texttt{resolution} (χρονική ανάλυση, συνήθως 15 λεπτά ή 1 ώρα). Τα δεδομένα αυτά χρησιμοποιούνται για την εξαγωγή τεχνικών παραμέτρων (ρυθμοί μεταβολής, ελάχιστη ισχύς, χρόνοι λειτουργίας/αδράνειας) μέσω στατιστικής ανάλυσης, όπως περιγράφεται στην Ενότητα~\ref{subsec:parameter-inference}.

\paragraph{Πίνακας unavailability\_of\_production\_and\_generation\_units}
Ο πίνακας αυτός περιέχει πληροφορίες διακοπών με τα πεδία: \texttt{asset\_code} (αντιστοιχεί στο \texttt{generation\_unit\_code}), \texttt{start\_outage\_utc} και \texttt{end\_outage\_utc} (περίοδος διακοπής), \texttt{type} (``Planned'' ή ``Forced''), \texttt{status} (``Active'', ``Cancelled'', ``Withdrawn''), και \texttt{available\_capacity\_mw} (εναπομένουσα διαθέσιμη ισχύς).

\subsection{Σχήμα Προσομοιώσεων}
\label{subsec:simulations-schema}

Το σχήμα \texttt{simulations} περιέχει τους πίνακες για την αποθήκευση των αποτελεσμάτων των προσομοιώσεων. Οι κύριοι πίνακες παρουσιάζονται στον Πίνακα~\ref{tab:simulations-tables}.

\begin{table}[htbp]
\centering
\caption{Κύριοι πίνακες του σχήματος \texttt{simulations}}
\label{tab:simulations-tables}
\begin{tabular}{@{}p{5cm}p{8cm}@{}}
\toprule
\textbf{Πίνακας} & \textbf{Περιγραφή} \\
\midrule
\texttt{energy\_prices} & Υπολογισμένες τιμές εκκαθάρισης \\
\texttt{optimization\_runs} & Μεταδεδομένα εκτελέσεων βελτιστοποίησης \\
\texttt{transmission\_flows} & Διασυνοριακές ροές ισχύος \\
\texttt{generator\_inferred\_parameters} & Εξαχθείσες τεχνικές παράμετροι μονάδων \\
\texttt{uc\_results} & Αποτελέσματα Unit Commitment (σύνοψη) \\
\texttt{uc\_generation} & Παραγωγή/δέσμευση ανά μονάδα και περίοδο \\
\texttt{uc\_net\_demand} & Καθαρή ζήτηση ανά περίοδο \\
\bottomrule
\end{tabular}
\end{table}

\paragraph{Πίνακας energy\_prices}
Ο πίνακας αυτός αποθηκεύει τις υπολογισμένες τιμές εκκαθάρισης με τα πεδία: \texttt{bidding\_zone} (ζώνη προσφοράς), \texttt{date} (ημερομηνία), \texttt{hour} (ώρα 1-24), \texttt{price} (τιμή σε €/MWh), \texttt{created\_at} (χρονοσφραγίδα δημιουργίας), και \texttt{run\_id} (αναφορά στην εκτέλεση βελτιστοποίησης).

\paragraph{Πίνακας transmission\_flows}
Ο πίνακας αυτός αποθηκεύει τις υπολογισμένες διασυνοριακές ροές με τα πεδία: \texttt{source\_zone} και \texttt{target\_zone} (ζώνες προέλευσης και προορισμού), \texttt{date} (ημερομηνία), \texttt{hour} (ώρα), \texttt{flow\_mw} (ροή σε MW), και \texttt{run\_id}.

\paragraph{Πίνακας generator\_inferred\_parameters}
Ο πίνακας αυτός λειτουργεί ως κρυφή μνήμη (cache) για τις τεχνικές παραμέτρους μονάδων παραγωγής που εξάγονται από ιστορικά δεδομένα. Τα βασικά πεδία περιλαμβάνουν: \texttt{generator\_code} (κωδικός μονάδας), \texttt{ramp\_up} και \texttt{ramp\_down} (ρυθμοί μεταβολής ως κλάσμα της $P^{\max}$ ανά ώρα), \texttt{p\_min} (ελάχιστη ισχύς ως κλάσμα της $P^{\max}$), \texttt{min\_uptime} και \texttt{min\_downtime} (ελάχιστοι χρόνοι λειτουργίας/αδράνειας σε ώρες), και \texttt{inference\_date} (ημερομηνία εξαγωγής για έλεγχο λήξης). Η χρήση του πίνακα αυτού μειώνει τον χρόνο ανάκτησης παραμέτρων από περίπου 17 λεπτά σε 2 δευτερόλεπτα.

\paragraph{Πίνακες uc\_results, uc\_generation, uc\_net\_demand}
Οι τρεις αυτοί πίνακες αποτελούν τον μηχανισμό κρυφής μνήμης (caching) για τα αποτελέσματα του επιλυτή Unit Commitment. Ο πίνακας \texttt{uc\_results} αποθηκεύει τη σύνοψη κάθε εκτέλεσης με κλειδί \texttt{(bidding\_zone, market\_date, code\_version)} και περιλαμβάνει τα πεδία: \texttt{status} (κατάσταση τερματισμού), \texttt{solver} (επιλυτής), \texttt{resolution\_minutes} (χρονική ανάλυση), \texttt{total\_cost}, \texttt{production\_cost}, \texttt{startup\_cost}, \texttt{noload\_cost} (συνιστώσες κόστους), και αριθμούς εκκινήσεων (\texttt{hot/warm/cold\_startups}). Ο πίνακας \texttt{uc\_generation} αποθηκεύει τα αναλυτικά αποτελέσματα παραγωγής, δέσμευσης και εκκίνησης ανά μονάδα και περίοδο, ενώ ο πίνακας \texttt{uc\_net\_demand} αποθηκεύει την καθαρή ζήτηση και την παραγωγή ΑΠΕ ανά περίοδο. Η αποθήκευση ανέρχεται σε περίπου 195 KB ανά ζώνη/ημέρα.


% =============================================================================
\section{Οικοσύστημα Julia/JuMP}
\label{sec:julia-jump}

Για την υλοποίηση των αλγορίθμων βελτιστοποίησης επιλέχθηκε η γλώσσα προγραμματισμού Julia σε συνδυασμό με το πλαίσιο μοντελοποίησης JuMP~\cite{bezanson2017, jump2017}. Αυτή η επιλογή βασίζεται σε τεχνικά πλεονεκτήματα που καθιστούν το συγκεκριμένο οικοσύστημα ιδανικό για επιστημονικούς υπολογισμούς και προβλήματα βελτιστοποίησης.

\subsection{Γλώσσα Προγραμματισμού Julia}
\label{subsec:julia-language}

Η Julia είναι μια σύγχρονη γλώσσα προγραμματισμού υψηλού επιπέδου σχεδιασμένη για επιστημονικούς υπολογισμούς. Τα κύρια χαρακτηριστικά της που την καθιστούν κατάλληλη για την παρούσα εφαρμογή είναι τα ακόλουθα.

Η Julia επιτυγχάνει υψηλή απόδοση μέσω της μεταγλώττισης just-in-time (JIT) με χρήση του LLVM compiler framework. Ο κώδικας Julia μεταγλωττίζεται σε εγγενή κώδικα μηχανής, επιτυγχάνοντας ταχύτητες συγκρίσιμες με C και Fortran, χωρίς την ανάγκη για χειροκίνητη βελτιστοποίηση ή κλήση εξωτερικών βιβλιοθηκών.

Παρά την υψηλή απόδοση, η σύνταξη της Julia είναι εκφραστική και φιλική στον χρήστη, παρόμοια με Python ή MATLAB. Αυτό επιτρέπει την ταχεία ανάπτυξη πρωτοτύπων και την εύκολη ανάγνωση του κώδικα.

Η Julia διαθέτει ένα εξελιγμένο σύστημα τύπων με πολλαπλή αποστολή (multiple dispatch) που επιτρέπει τη δημιουργία γενικευμένου κώδικα που είναι ταυτόχρονα αποτελεσματικός και επαναχρησιμοποιήσιμος.

Η Julia παρέχει ενσωματωμένη υποστήριξη για παραλληλισμό και κατανεμημένους υπολογισμούς, διευκολύνοντας την κλιμάκωση σε πολλαπλούς πυρήνες ή μηχανές.

\subsection{Πλαίσιο Μοντελοποίησης JuMP}
\label{subsec:jump-framework}

Το JuMP (Julia for Mathematical Programming) είναι ένα domain-specific language (DSL) ενσωματωμένο στη Julia για τη διατύπωση και επίλυση προβλημάτων μαθηματικού προγραμματισμού. Αποτελεί ένα από τα πιο δημοφιλή και ώριμα πλαίσια μοντελοποίησης βελτιστοποίησης στον ακαδημαϊκό και βιομηχανικό χώρο~\cite{jump2017}.

Τα πλεονεκτήματα του JuMP περιλαμβάνουν μια αλγεβρική μοντελοποίηση που επιτρέπει τη διατύπωση προβλημάτων βελτιστοποίησης με τρόπο που μοιάζει με τη μαθηματική τους διατύπωση. Η σύνταξη για μεταβλητές, περιορισμούς και αντικειμενικές συναρτήσεις είναι διαισθητική και εκφραστική. Το JuMP υποστηρίζει πολλαπλούς επιλυτές (solvers) μέσω μιας ενιαίας διεπαφής, επιτρέποντας την εύκολη εναλλαγή μεταξύ επιλυτών χωρίς αλλαγές στον κώδικα μοντελοποίησης. Η απόδοση του JuMP είναι συγκρίσιμη ή καλύτερη από εμπορικά εργαλεία μοντελοποίησης όπως το AMPL και το GAMS. Τέλος, το JuMP είναι εντελώς ελεύθερο και ανοιχτού κώδικα, διατίθεται υπό την άδεια MIT.

Ένα τυπικό παράδειγμα χρήσης του JuMP για τη δημιουργία ενός μοντέλου βελτιστοποίησης φαίνεται στο ακόλουθο απόσπασμα κώδικα:

\begin{samepage}
\begin{minted}[fontsize=\small,linenos,frame=lines]{julia}
using JuMP, HiGHS

# Δημιουργία μοντέλου με τον επιλυτή HiGHS
model = Model(HiGHS.Optimizer)

# Ορισμός μεταβλητών
@variable(model, x >= 0)
@variable(model, y >= 0)

# Ορισμός περιορισμών
@constraint(model, x + 2y <= 10)
@constraint(model, 3x + y <= 15)

# Ορισμός αντικειμενικής συνάρτησης
@objective(model, Max, 2x + 3y)

# Επίλυση
optimize!(model)
\end{minted}
\end{samepage}

\subsection{Επιλυτής HiGHS}
\label{subsec:highs-solver}

Ως κύριος επιλυτής για τα προβλήματα βελτιστοποίησης χρησιμοποιείται ο HiGHS, ένας σύγχρονος επιλυτής ανοιχτού κώδικα για προβλήματα γραμμικού προγραμματισμού (LP), μικτού ακέραιου προγραμματισμού (MIP) και τετραγωνικού προγραμματισμού (QP)~\cite{highs2024}.

Ο HiGHS αναπτύσσεται στο Πανεπιστήμιο του Εδιμβούργου και αποτελεί έναν από τους ταχύτερους ανοιχτού κώδικα επιλυτές. Περιλαμβάνει υλοποιήσεις του αλγορίθμου dual revised simplex με εκτεταμένη παραλληλοποίηση, μεθόδων εσωτερικού σημείου (interior point methods), και αλγορίθμων branch-and-bound για MIP με cutting planes και heuristics.

Η διεπαφή με τη Julia παρέχεται μέσω του πακέτου HiGHS.jl, το οποίο ενσωματώνεται απρόσκοπτα με το JuMP. Η επιλογή του HiGHS ως προεπιλεγμένου επιλυτή βασίζεται στην άδεια ανοιχτού κώδικα (MIT) που επιτρέπει ελεύθερη χρήση, στην υψηλή απόδοση συγκρίσιμη με εμπορικούς επιλυτές, στην ενεργή ανάπτυξη και υποστήριξη, και στη συμβατότητα με όλα τα απαιτούμενα είδη προβλημάτων (LP, MILP).

Το σύστημα υποστηρίζει επίσης εμπορικούς επιλυτές όπως ο Gurobi και ο CPLEX για περιπτώσεις που απαιτείται ακόμη υψηλότερη απόδοση, με την εναλλαγή να γίνεται μέσω μιας απλής παραμέτρου.

\subsection{Βασικά Πακέτα Julia}
\label{subsec:julia-packages}

Εκτός από το JuMP και τον HiGHS, το σύστημα χρησιμοποιεί μια σειρά από πακέτα Julia για διάφορες λειτουργίες. Το πακέτο DataFrames.jl χρησιμοποιείται για τη διαχείριση και επεξεργασία δεδομένων σε μορφή πινάκων, με λειτουργικότητα παρόμοια με τα pandas της Python. Το πακέτο LibPQ.jl παρέχει σύνδεση με τη βάση δεδομένων PostgreSQL μέσω της βιβλιοθήκης libpq, υποστηρίζοντας παραμετροποιημένα ερωτήματα και διαχείριση συνδέσεων. Το πακέτο CSV.jl χρησιμοποιείται για την ανάγνωση και εγγραφή αρχείων CSV με υψηλή απόδοση. Το πακέτο Dates.jl, που αποτελεί μέρος της standard library, χρησιμοποιείται για τη διαχείριση ημερομηνιών και χρόνων. Το πακέτο Statistics.jl, επίσης μέρος της standard library, χρησιμοποιείται για στατιστικούς υπολογισμούς---κυρίως τη συνάρτηση \texttt{quantile} για τον υπολογισμό εκατοστημορίων κατά την εξαγωγή τεχνικών παραμέτρων μονάδων από ιστορικά δεδομένα. Τέλος, το πακέτο DotEnv.jl χρησιμοποιείται για τη φόρτωση μεταβλητών περιβάλλοντος από αρχεία \texttt{.env} για τη διαχείριση διαπιστευτηρίων.


% =============================================================================
\section{Αρχιτεκτονική Εφαρμογής}
\label{sec:application-architecture}

Η εφαρμογή Julia οργανώνεται ως ένα ενιαίο module ονόματι \texttt{Euphemia} που περιέχει όλη τη λειτουργικότητα του συστήματος. Η δομή του module ακολουθεί τις βέλτιστες πρακτικές της Julia για οργάνωση κώδικα και επαναχρησιμοποίηση~\cite{julia-docs}.

\subsection{Δομή Module}
\label{subsec:module-structure}

Το module \texttt{Euphemia} αποτελείται από τα ακόλουθα υπο-modules και αρχεία, όπως φαίνεται στο Σχήμα~\ref{fig:module-structure}:

\begin{figure}[htbp]
\centering
\begin{verbatim}
Euphemia.jl (Main Module)
├── dbutils.jl                 # Database utilities
├── Generators.jl              # Generator data model
├── Loads.jl                   # Load data model
├── Renewables.jl              # RES forecasts
├── FuelTypeParameters.jl      # Fuel-specific parameters
├── MarketOrders.jl            # Order types
├── UnitCommitment.jl          # UC solver
├── BiddingStrategy.jl         # UC → Market orders
├── Network.jl                 # ATC constraints
├── MPCC.jl                    # Market clearing
└── AlternativeOrderBook.jl    # Fast order generation
\end{verbatim}
\caption{Δομή του module \texttt{Euphemia}.}
\label{fig:module-structure}
\end{figure}

\paragraph{Κεντρικό Module (Euphemia.jl)}
Αποτελεί το σημείο εισόδου της εφαρμογής. Διαχειρίζεται τις εξαρτήσεις, συντονίζει τα υπο-modules, και εξάγει τη δημόσια διεπαφή (public API). Περιλαμβάνει επίσης τη λογική επιλογής επιλυτή και την κεντρική ροή εργασιών.

\paragraph{Βοηθητικά Βάσης Δεδομένων (dbutils.jl)}
Περιέχει συναρτήσεις για τη σύνδεση με τη βάση δεδομένων PostgreSQL και την εκτέλεση ερωτημάτων. Υλοποιεί το pattern \texttt{withdb} για ασφαλή διαχείριση συνδέσεων με αυτόματο κλείσιμο.

\paragraph{Μοντέλα Δεδομένων (Generators.jl, Loads.jl, Renewables.jl)}
Ορίζουν τις δομές δεδομένων για τις οντότητες του συστήματος και παρέχουν συναρτήσεις για την ανάκτησή τους από τη βάση δεδομένων. Η δομή \texttt{Generator} περιλαμβάνει πεδία για τον κωδικό, το όνομα, τη ζώνη προσφοράς, τον τύπο καυσίμου, τα όρια ισχύος, το οριακό κόστος, και προαιρετικούς τεχνικούς περιορισμούς (ρυθμοί μεταβολής, ελάχιστοι χρόνοι λειτουργίας/αδράνειας) που εξάγονται αυτόματα από ιστορικά δεδομένα. Η δομή \texttt{InitialConditions} αναπαριστά την αρχική κατάσταση κάθε μονάδας στη χρονική στιγμή $t=0$ (δέσμευση, ισχύς εξόδου, θερμοκρασιακή κατάσταση) και εξάγεται από ιστορικά δεδομένα 72 ωρών μέσω μαζικού ερωτήματος. Το αρχείο \texttt{Generators.jl} περιλαμβάνει τις συναρτήσεις εξαγωγής παραμέτρων (\texttt{infer\_ramp\_rates}, \texttt{infer\_p\_min}, \texttt{infer\_uptime\_downtime}), τη συνάρτηση \texttt{get\_initial\_conditions} για αυτόματο προσδιορισμό αρχικών συνθηκών, και τον μηχανισμό αποθήκευσης σε κρυφή μνήμη βάσης δεδομένων.

\paragraph{Τύποι Εντολών (MarketOrders.jl)}
Ορίζει τις δομές δεδομένων για τους διάφορους τύπους εντολών αγοράς, συμπεριλαμβανομένων των \texttt{SimpleOrder}, \texttt{BlockOrder}, \texttt{LinkedBlockOrder} κ.λπ.

\paragraph{Unit Commitment (UnitCommitment.jl)}
Υλοποιεί τον επιλυτή για το πρόβλημα Unit Commitment χρησιμοποιώντας JuMP/HiGHS. Περιλαμβάνει τη διατύπωση του μοντέλου με τριμερή αντικειμενική συνάρτηση (κόστος παραγωγής, κόστος εκκίνησης ανά θερμοκρασιακή κατάσταση, κόστος αδράνειας), την ενσωμάτωση αρχικών συνθηκών, ρυθμιζόμενες παραμέτρους επιλυτή (MIP gap, χρονικό όριο μέσω MOI), καθώς και μηχανισμό αποθήκευσης αποτελεσμάτων σε κρυφή μνήμη βάσης δεδομένων.

\paragraph{Στρατηγική Προσφορών (BiddingStrategy.jl)}
Υλοποιεί τη μετατροπή των αποτελεσμάτων του Unit Commitment σε εντολές αγοράς σύμφωνα με τις τρεις στρατηγικές που περιγράφηκαν στο Κεφάλαιο~\ref{sec:bidding-strategies}.

\paragraph{Δίκτυο (Network.jl)}
Διαχειρίζεται τους περιορισμούς δικτύου και τη διαθέσιμη ικανότητα μεταφοράς (ATC) μεταξύ ζωνών προσφοράς. Παρέχει λειτουργίες για την ανάκτηση δεδομένων ATC από το ENTSO-E και την προσθήκη περιορισμών στο μοντέλο βελτιστοποίησης.

\paragraph{Εκκαθάριση Αγοράς (MPCC.jl)}
Υλοποιεί τον αλγόριθμο εκκαθάρισης της αγοράς με τη μέθοδο MPCC. Περιλαμβάνει τη δομή \texttt{MPCCOrderBook} για τη διαχείριση των εντολών και τη συνάρτηση \texttt{solve\_mpcc\_market\_clearing} για την επίλυση.

\subsection{Ροή Εργασιών}
\label{subsec:workflow}

Η κεντρική ροή εργασιών του συστήματος για τον υπολογισμό τιμών αγοράς περιλαμβάνει τα ακόλουθα βήματα. Αρχικά, γίνεται φόρτωση δεδομένων από τη βάση, συμπεριλαμβανομένων των διαθέσιμων μονάδων παραγωγής (με φιλτράρισμα μη διαθεσιμοτήτων), της ζήτησης για την επιλεγμένη ημερομηνία, και των προβλέψεων παραγωγής ΑΠΕ. Στη συνέχεια, επιλύεται το πρόβλημα Unit Commitment για τον προσδιορισμό του βέλτιστου προγράμματος λειτουργίας. Ακολουθεί η μετατροπή της λύσης του UC σε εντολές αγοράς σύμφωνα με την επιλεγμένη στρατηγική. Εκτελείται η εκκαθάριση της αγοράς μέσω του αλγορίθμου MPCC, που καθορίζει τις τιμές ισορροπίας. Τέλος, τα αποτελέσματα αποθηκεύονται στη βάση δεδομένων.

Η κύρια συνάρτηση για την εκτέλεση αυτής της ροής είναι η \texttt{generate\_energy\_prices}:

\begin{samepage}
\begin{minted}[fontsize=\small,linenos,frame=lines]{julia}
# Παράδειγμα χρήσης για μονοζωνική εκκαθάριση
prices = generate_energy_prices("GR", Date(2024, 6, 15);
    order_method = :uc_based,
    markup_factor = 1.1,
    save_to_db = true
)

# Πολυζωνική εκκαθάριση με περιορισμούς ATC
result = run_multi_zone_market_clearing(Date(2024, 6, 15);
    zones = ["GR", "BG", "RO"],
    order_method = :alternative,
    save_to_db = true
)
\end{minted}
\end{samepage}

\subsection{Διαχείριση Επιλυτών}
\label{subsec:solver-management}

Το σύστημα υποστηρίζει πολλαπλούς επιλυτές μέσω ενός μηχανισμού αυτόματης επιλογής και caching. Η συνάρτηση \texttt{select\_solver} ελέγχει τη διαθεσιμότητα των εγκατεστημένων επιλυτών (HiGHS, Gurobi, CPLEX) και επιλέγει τον κατάλληλο βάσει προτίμησης ή διαθεσιμότητας.

Για τον επιλυτή Gurobi, ο οποίος απαιτεί άδεια χρήσης με σημαντικό κόστος αρχικοποίησης, υλοποιήθηκε ένας μηχανισμός caching του περιβάλλοντος (environment) που αποφεύγει την επανάληψη της διαδικασίας αυθεντικοποίησης σε διαδοχικές κλήσεις.

\begin{samepage}
\begin{minted}[fontsize=\small,linenos,frame=lines]{julia}
# Αυτόματη επιλογή επιλυτή
optimizer, solver_name = select_solver("auto")

# Προτίμηση συγκεκριμένου επιλυτή
optimizer, solver_name = select_solver("gurobi")

# Δημιουργία μοντέλου με τον επιλεγμένο επιλυτή
model = Model(optimizer)
\end{minted}
\end{samepage}
